\chapter{La gestione avanzata dei processi}
\label{cha:proc_advanced}

In questo capitolo affronteremo gli argomenti relativi alla gestione avanzata
dei processi. Inizieremo con le funzioni che attengono alla gestione avanzata
della sicurezza, passando poi a quelle relative all'analisi ed al controllo
dell'esecuzione, e alle funzioni per le modalità avanzate di creazione dei
processi e l'uso dei cosiddetti \textit{namespace}. Infine affronteremo le
\textit{sytem call} attinenti ad una serie di funzionalità specialistiche come
la gestione della virgola mobile, le porte di I/O ecc.

\section{La gestione avanzata della sicurezza}
\label{sec:process_security}

Tratteremo in questa sezione le funzionalità più avanzate relative alla
gestione della sicurezza ed il controllo degli accessi all'interno dei
processi, a partire dalle \textit{capabilities} e dalle funzionalità del
cosiddetto \textit{Secure Computing}. Esamineremo inoltre le altre
funzionalità relative alla sicurezza come la gestione delle chiavi
crittografiche e varie estensioni e funzionalità disponibili su questo
argomento.


\subsection{La gestione delle \textit{capabilities}}
\label{sec:proc_capabilities}

\itindbeg{capabilities} 

Come accennato in sez.~\ref{sec:proc_access_id} l'architettura classica della
gestione dei privilegi in un sistema unix-like ha il sostanziale problema di
fornire all'amministratore dei poteri troppo ampi, il che comporta che anche
quando si siano predisposte delle misure di protezione per in essere in grado
di difendersi dagli effetti di una eventuale compromissione del sistema (come
montare un filesystem in sola lettura per impedirne modifiche, o marcare un
file come immutabile) una volta che questa sia stata effettuata e si siano
ottenuti i privilegi di amministratore, queste misure potranno essere comunque
rimosse (nei casi elencati nella precedente nota si potrà sempre rimontare il
sistema in lettura-scrittura, o togliere l'attributo di immutabilità).

Il problema consiste nel fatto che nell'architettura tradizionale di un
sistema unix-like i controlli di accesso sono basati su un solo livello di
separazione: per i processi normali essi sono posti in atto, mentre per i
processi con i privilegi di amministratore essi non vengono neppure eseguiti.
Per questo motivo non era previsto alcun modo per evitare che un processo con
diritti di amministratore non potesse eseguire certe operazioni, o per cedere
definitivamente alcuni privilegi da un certo momento in poi. 

Per risolvere questo problema sono possibili varie soluzioni, ad esempio dai
kernel della serie 2.5 è stata introdotta la struttura dei
\itindex{Linux~Security~Modules~(LSM)} \textit{Linux Security Modules} che han
permesso di aggiungere varie forme di \itindex{Mandatory~Access~Control~(DAC)}
\textit{Mandatory Access Control} (MAC), in cui si potessero parcellizzare e
controllare nei minimi dettagli tutti i privilegi e le modalità in cui questi
possono essere usati dai programmi e trasferiti agli utenti, con la creazione
di varie estensioni (come \textit{SELinux}, \textit{Smack}, \textit{Tomoyo},
\textit{AppArmor}) che consentono di superare l'architettura tradizionale dei
permessi basati sul modello classico del controllo di accesso chiamato
\itindex{Discrectionary~Access~Control~(DAC)} \textit{Discrectionary Access
  Control} (DAC).

Ma già in precedenza, a partire dai kernel della serie 2.2, era stato
introdotto un meccanismo, detto \textit{capabilities}, per consentire di
suddividere i vari privilegi tradizionalmente associati all'amministratore in
un insieme di \textsl{capacità} distinte.  L'idea era che queste capacità
potessero essere abilitate e disabilitate in maniera indipendente per ciascun
processo con privilegi di amministratore, permettendo così una granularità
molto più fine nella distribuzione degli stessi, che evitasse la situazione
originaria di ``\textsl{tutto o nulla}''.

\itindbeg{file~capabilities}

Il meccanismo completo delle \textit{capabilities} (l'implementazione si rifà
ad una bozza di quello che doveva diventare lo standard POSIX.1e, poi
abbandonato) prevede inoltre la possibilità di associare le stesse ai singoli
file eseguibili, in modo da poter stabilire quali capacità possono essere
utilizzate quando viene messo in esecuzione uno specifico programma; ma il
supporto per questa funzionalità, chiamata \textit{file capabilities}, è stato
introdotto soltanto a partire dal kernel 2.6.24. Fino ad allora doveva essere
il programma stesso ad eseguire una riduzione esplicita delle sue capacità,
cosa che ha reso l'uso di questa funzionalità poco diffuso, vista la presenza
di meccanismi alternativi per ottenere limitazioni delle capacità
dell'amministratore a livello di sistema operativo, come \textit{SELinux}.

Con questo supporto e con le ulteriori modifiche introdotte con il kernel
2.6.25 il meccanismo delle \textit{capabilities} è stato totalmente
rivoluzionato, rendendolo più aderente alle intenzioni originali dello
standard POSIX, rimuovendo il significato che fino ad allora aveva avuto la
capacità \const{CAP\_SETPCAP}, e cambiando le modalità di funzionamento del
cosiddetto \textit{capabilities bounding set}. Ulteriori modifiche sono state
apportate con il kernel 2.6.26 per consentire la rimozione non ripristinabile
dei privilegi di amministratore. Questo fa sì che il significato ed il
comportamento del kernel finisca per dipendere dalla versione dello stesso e
dal fatto che le nuove \textit{file capabilities} siano abilitate o meno. Per
capire meglio la situazione e cosa è cambiato conviene allora spiegare con
maggiori dettagli come funziona il meccanismo delle \textit{capabilities}.

Il primo passo per frazionare i privilegi garantiti all'amministratore,
supportato fin dalla introduzione iniziale del kernel 2.2, è stato quello in
cui a ciascun processo sono stati associati tre distinti insiemi di
\textit{capabilities}, denominati rispettivamente \textit{permitted},
\textit{inheritable} ed \textit{effective}. Questi insiemi vengono mantenuti
in forma di tre diverse maschere binarie,\footnote{il kernel li mantiene, come
  i vari identificatori di sez.~\ref{sec:proc_setuid}, all'interno della
  \texttt{task\_struct} di ciascun processo (vedi
  fig.~\ref{fig:proc_task_struct}), nei tre campi \texttt{cap\_effective},
  \texttt{cap\_inheritable}, \texttt{cap\_permitted} del tipo
  \texttt{kernel\_cap\_t}; questo era, fino al kernel 2.6.25 definito come
  intero a 32 bit per un massimo di 32 \textit{capabilities} distinte,
  attualmente è stato aggiornato ad un vettore in grado di mantenerne fino a
  64.} in cui ciascun bit corrisponde ad una capacità diversa.

L'utilizzo di tre distinti insiemi serve a fornire una interfaccia flessibile
per l'uso delle \textit{capabilities}, con scopi analoghi a quelli per cui
sono mantenuti i diversi insiemi di identificatori di
sez.~\ref{sec:proc_setuid}; il loro significato, che è rimasto sostanzialmente
lo stesso anche dopo le modifiche seguite alla introduzione delle
\textit{file capabilities} è il seguente:
\begin{basedescript}{\desclabelwidth{2.1cm}\desclabelstyle{\nextlinelabel}}
\item[\textit{permitted}] l'insieme delle \textit{capabilities}
  ``\textsl{permesse}'', cioè l'insieme di quelle capacità che un processo
  \textsl{può} impostare come \textsl{effettive} o come
  \textsl{ereditabili}. Se un processo cancella una capacità da questo insieme
  non potrà più riassumerla.\footnote{questo nei casi ordinari, sono
    previste però una serie di eccezioni, dipendenti anche dal tipo di
    supporto, che vedremo meglio in seguito dato il notevole intreccio nella
    casistica.}
\item[\textit{inheritable}] l'insieme delle \textit{capabilities}
  ``\textsl{ereditabili}'', cioè di quelle che verranno trasmesse come insieme
  delle \textsl{permesse} ad un nuovo programma eseguito attraverso una
  chiamata ad \func{exec}.
\item[\textit{effective}] l'insieme delle \textit{capabilities}
  ``\textsl{effettive}'', cioè di quelle che vengono effettivamente usate dal
  kernel quando deve eseguire il controllo di accesso per le varie operazioni
  compiute dal processo.
\label{sec:capabilities_set}
\end{basedescript}

Con l'introduzione delle \textit{file capabilities} sono stati introdotti
altri tre insiemi associabili a ciascun file.\footnote{la realizzazione viene
  eseguita con l'uso di uno specifico attributo esteso,
  \texttt{security.capability}, la cui modifica è riservata, (come illustrato
  in sez.~\ref{sec:file_xattr}) ai processi dotato della capacità
  \const{CAP\_SYS\_ADMIN}.} Le \textit{file capabilities} hanno effetto
soltanto quando il file che le porta viene eseguito come programma con una
\func{exec}, e forniscono un meccanismo che consente l'esecuzione dello stesso
con maggiori privilegi; in sostanza sono una sorta di estensione del
\acr{suid} bit limitato ai privilegi di amministratore. Anche questi tre
insiemi sono identificati con gli stessi nomi di quello dei processi, ma il
loro significato è diverso:
\begin{basedescript}{\desclabelwidth{2.1cm}\desclabelstyle{\nextlinelabel}}
\item[\textit{permitted}] (chiamato originariamente \textit{forced}) l'insieme
  delle capacità che con l'esecuzione del programma verranno aggiunte alle
  capacità \textsl{permesse} del processo.
\item[\textit{inheritable}] (chiamato originariamente \textit{allowed})
  l'insieme delle capacità che con l'esecuzione del programma possono essere
  ereditate dal processo originario (che cioè non vengono tolte
  dall'\textit{inheritable set} del processo originale all'esecuzione di
  \func{exec}).
\item[\textit{effective}] in questo caso non si tratta di un insieme ma di un
  unico valore logico; se attivo all'esecuzione del programma tutte le
  capacità che risulterebbero \textsl{permesse} verranno pure attivate,
  inserendole automaticamente nelle \textsl{effettive}, se disattivato nessuna
  capacità verrà attivata (cioè l'\textit{effective set} resterà vuoto).
\end{basedescript}

\itindbeg{capabilities~bounding~set}

Infine come accennato, esiste un ulteriore insieme, chiamato
\textit{capabilities bounding set}, il cui scopo è quello di costituire un
limite alle capacità che possono essere attivate per un programma. Il suo
funzionamento però è stato notevolmente modificato con l'introduzione delle
\textit{file capabilities} e si deve pertanto prendere in considerazione una
casistica assai complessa.

Per i kernel fino al 2.6.25, o se non si attiva il supporto per le
\textit{file capabilities}, il \textit{capabilities bounding set} è un
parametro generale di sistema, il cui valore viene riportato nel file
\sysctlfiled{kernel/cap-bound}. Il suo valore iniziale è definito in sede di
compilazione del kernel, e da sempre ha previsto come default la presenza di
tutte le \textit{capabilities} eccetto \const{CAP\_SETPCAP}. In questa
situazione solo il primo processo eseguito nel sistema (quello con
\textsl{pid} 1, di norma \texttt{/sbin/init}) ha la possibilità di
modificarlo; ogni processo eseguito successivamente, se dotato dei privilegi
di amministratore, è in grado soltanto di rimuovere una delle
\textit{capabilities} già presenti dell'insieme.\footnote{per essere precisi
  occorre la capacità \const{CAP\_SYS\_MODULE}.}

In questo caso l'effetto complessivo del \textit{capabilities bounding set} è
che solo le capacità in esso presenti possono essere trasmesse ad un altro
programma attraverso una \func{exec}. Questo in sostanza significa che se un
qualunque programma elimina da esso una capacità, considerato che
\texttt{init} (almeno nelle versioni ordinarie) non supporta la reimpostazione
del \textit{bounding set}, questa non sarà più disponibile per nessun processo
a meno di un riavvio, eliminando così in forma definitiva quella capacità per
tutti, compreso l'amministratore.\footnote{la qual cosa, visto il default
  usato per il \textit{capabilities bounding set}, significa anche che
  \const{CAP\_SETPCAP} non è stata praticamente mai usata nella sua forma
  originale.}

Con il kernel 2.6.25 e le \textit{file capabilities} il \textit{bounding set}
è diventato una proprietà di ciascun processo, che viene propagata invariata
sia attraverso una \func{fork} che una \func{exec}. In questo caso il file
\sysctlfile{kernel/cap-bound} non esiste e \texttt{init} non ha nessun
ruolo speciale, inoltre in questo caso all'avvio il valore iniziale prevede la
presenza di tutte le capacità (compresa \const{CAP\_SETPCAP}). 

Con questo nuovo meccanismo il \textit{bounding set} continua a ricoprire un
ruolo analogo al precedente nel passaggio attraverso una \func{exec}, come
limite alle capacità che possono essere aggiunte al processo in quanto
presenti nel \textit{permitted set} del programma messo in esecuzione, in
sostanza il nuovo programma eseguito potrà ricevere una capacità presente nel
suo \textit{permitted set} (quello del file) solo se questa è anche nel
\textit{bounding set} (del processo). In questo modo si possono rimuovere
definitivamente certe capacità da un processo, anche qualora questo dovesse
eseguire un programma privilegiato che prevede di riassegnarle.

Si tenga presente però che in questo caso il \textit{bounding set} blocca
esclusivamente le capacità indicate nel \textit{permitted set} del programma
che verrebbero attivate in caso di esecuzione, e non quelle eventualmente già
presenti nell'\textit{inheritable set} del processo (ad esempio perché
presenti prima di averle rimosse dal \textit{bounding set}). In questo caso
eseguendo un programma che abbia anche lui dette capacità nel suo
\textit{inheritable set} queste verrebbero assegnate.

In questa seconda versione inoltre il \textit{bounding set} costituisce anche
un limite per le capacità che possono essere aggiunte all'\textit{inheritable
  set} del processo stesso con \func{capset}, sempre nel senso che queste
devono essere presenti nel \textit{bounding set} oltre che nel
\textit{permitted set} del processo. Questo limite vale anche per processi con
i privilegi di amministratore,\footnote{si tratta sempre di avere la
  \textit{capability} \const{CAP\_SETPCAP}.} per i quali invece non vale la
condizione che le \textit{capabilities} da aggiungere nell'\textit{inheritable
  set} debbano essere presenti nel proprio \textit{permitted set}.\footnote{lo
  scopo anche in questo caso è ottenere una rimozione definitiva della
  possibilità di passare una capacità rimossa dal \textit{bounding set}.}

Come si può notare per fare ricorso alle \textit{capabilities} occorre
comunque farsi carico di una notevole complessità di gestione, aggravata dalla
presenza di una radicale modifica del loro funzionamento con l'introduzione
delle \textit{file capabilities}. Considerato che il meccanismo originale era
incompleto e decisamente problematico nel caso di programmi che non ne
sapessero tener conto,\footnote{il problema di sicurezza originante da questa
  caratteristica venne alla ribalta con \texttt{sendmail}, in cui, riuscendo a
  rimuovere \const{CAP\_SETGID} dall'\textit{inheritable set} di un processo,
  si ottenne di far fallire \func{setuid} in maniera inaspettata per il
  programma (che aspettandosi sempre il successo della funzione non ne
  controllava lo stato di uscita) con la conseguenza di fargli fare come
  amministratore operazioni che altrimenti sarebbero state eseguite, senza
  poter apportare danni, da utente normale.}  ci soffermeremo solo sulla
implementazione completa presente a partire dal kernel 2.6.25, tralasciando
ulteriori dettagli riguardo la versione precedente.

Riassumendo le regole finora illustrate tutte le \textit{capabilities} vengono
ereditate senza modifiche attraverso una \func{fork} mentre, indicati con
\texttt{orig\_*} i valori degli insiemi del processo chiamante, con
\texttt{file\_*} quelli del file eseguito e con \texttt{bound\_set} il
\textit{capabilities bounding set}, dopo l'invocazione di \func{exec} il
processo otterrà dei nuovi insiemi di capacità \texttt{new\_*} secondo la
formula espressa dal seguente pseudo-codice C:

\includecodesnip{listati/cap-results.c}

% \begin{figure}[!htbp]
%   \footnotesize \centering
%   \begin{minipage}[c]{12cm}
%     \includecodesnip{listati/cap-results.c}
%   \end{minipage}
%   \caption{Espressione della modifica delle \textit{capabilities} attraverso
%     una \func{exec}.}
%   \label{fig:cap_across_exec}
% \end{figure}

\noindent e si noti come in particolare il \textit{capabilities bounding set}
non venga comunque modificato e resti lo stesso sia attraverso una \func{fork}
che attraverso una \func{exec}.


\itindend{capabilities~bounding~set}

A queste regole se ne aggiungono delle altre che servono a riprodurre il
comportamento tradizionale di un sistema unix-like in tutta una serie di
circostanze. La prima di queste è relativa a quello che avviene quando si
esegue un file senza \textit{capabilities}; se infatti si considerasse questo
equivalente al non averne assegnata alcuna, non essendo presenti capacità né
nel \textit{permitted set} né nell'\textit{inheritable set} del file,
nell'esecuzione di un qualunque programma l'amministratore perderebbe tutti i
privilegi originali dal processo.

Per questo motivo se un programma senza \textit{capabilities} assegnate viene
eseguito da un processo con \ids{UID} reale 0, esso verrà trattato come
se tanto il \textit{permitted set} che l'\textit{inheritable set} fossero con
tutte le \textit{capabilities} abilitate, con l'\textit{effective set} attivo,
col risultato di fornire comunque al processo tutte le capacità presenti nel
proprio \textit{bounding set}. Lo stesso avviene quando l'eseguibile ha attivo
il \acr{suid} bit ed appartiene all'amministratore, in entrambi i casi si
riesce così a riottenere il comportamento classico di un sistema unix-like.

Una seconda circostanza è quella relativa a cosa succede alle
\textit{capabilities} di un processo nelle possibili transizioni da \ids{UID}
nullo a \ids{UID} non nullo o viceversa (corrispondenti rispettivamente a
cedere o riottenere i privilegi di amministratore) che si possono effettuare
con le varie funzioni viste in sez.~\ref{sec:proc_setuid}. In questo caso la
casistica è di nuovo alquanto complessa, considerata anche la presenza dei
diversi gruppi di identificatori illustrati in tab.~\ref{tab:proc_uid_gid}, si
avrà allora che:
\begin{enumerate*}
\item se si passa da \ids{UID} effettivo nullo a non nullo
  l'\textit{effective set} del processo viene totalmente azzerato, se
  viceversa si passa da \ids{UID} effettivo non nullo a nullo il
  \textit{permitted set} viene copiato nell'\textit{effective set};
\item se si passa da \textit{file system} \ids{UID} nullo a non nullo verranno
  cancellate dall'\textit{effective set} del processo tutte le capacità
  attinenti i file, e cioè \const{CAP\_LINUX\_IMMUTABLE}, \const{CAP\_MKNOD},
  \const{CAP\_DAC\_OVERRIDE}, \const{CAP\_DAC\_READ\_SEARCH},
  \const{CAP\_MAC\_OVERRIDE}, \const{CAP\_CHOWN}, \const{CAP\_FSETID} e
  \const{CAP\_FOWNER} (le prime due a partire dal kernel 2.2.30), nella
  transizione inversa verranno invece inserite nell'\textit{effective set}
  quelle capacità della precedente lista che sono presenti nel suo
  \textit{permitted set}.
\item se come risultato di una transizione riguardante gli identificativi dei
  gruppi \textit{real}, \textit{saved} ed \textit{effective} in cui si passa
  da una situazione in cui uno di questi era nullo ad una in cui sono tutti
  non nulli,\footnote{in sostanza questo è il caso di quando si chiama
    \func{setuid} per rimuovere definitivamente i privilegi di amministratore
    da un processo.} verranno azzerati completamente sia il \textit{permitted
    set} che l'\textit{effective set}.
\end{enumerate*}
\label{sec:capability-uid-transition}

La combinazione di tutte queste regole consente di riprodurre il comportamento
ordinario di un sistema di tipo Unix tradizionale, ma può risultare
problematica qualora si voglia passare ad una configurazione di sistema
totalmente basata sull'applicazione delle \textit{capabilities}; in tal caso
infatti basta ad esempio eseguire un programma con \acr{suid} bit di proprietà
dell'amministratore per far riottenere ad un processo tutte le capacità
presenti nel suo \textit{bounding set}, anche se si era avuta la cura di
cancellarle dal \textit{permitted set}.

\itindbeg{securebits}

Per questo motivo a partire dal kernel 2.6.26, se le \textit{file
  capabilities} sono abilitate, ad ogni processo viene stata associata una
ulteriore maschera binaria, chiamata \textit{securebits flags}, su cui sono
mantenuti una serie di flag (vedi tab.~\ref{tab:securebits_values}) il cui
valore consente di modificare queste regole speciali che si applicano ai
processi con \ids{UID} nullo. La maschera viene sempre mantenuta
attraverso una \func{fork}, mentre attraverso una \func{exec} viene sempre
cancellato il flag \const{SECURE\_KEEP\_CAPS}.

\begin{table}[htb]
  \centering
  \footnotesize
  \begin{tabular}{|l|p{10cm}|}
    \hline
    \textbf{Flag} & \textbf{Descrizione} \\
    \hline
    \hline
    \constd{SECURE\_KEEP\_CAPS}&Il processo non subisce la cancellazione delle
                                sue \textit{capabilities} quando tutti i suoi
                                \ids{UID} passano ad un valore non
                                nullo (regola di compatibilità per il cambio
                                di \ids{UID} n.~3 del precedente
                                elenco), sostituisce il precedente uso
                                dell'operazione \const{PR\_SET\_KEEPCAPS} di
                                \func{prctl}.\\
    \constd{SECURE\_NO\_SETUID\_FIXUP}&Il processo non subisce le modifiche
                                delle sue \textit{capabilities} nel passaggio
                                da nullo a non nullo degli \ids{UID}
                                dei gruppi \textit{effective} e
                                \textit{file system} (regole di compatibilità
                                per il cambio di \ids{UID} nn.~1 e 2 del
                                precedente elenco).\\
    \constd{SECURE\_NOROOT}   & Il processo non assume nessuna capacità
                                aggiuntiva quando esegue un programma, anche
                                se ha \ids{UID} nullo o il programma ha
                                il \acr{suid} bit attivo ed appartiene
                                all'amministratore (regola di compatibilità
                                per l'esecuzione di programmi senza
                                \textit{capabilities}).\\
    \hline
  \end{tabular}
  \caption{Costanti identificative dei flag che compongono la maschera dei
    \textit{securebits}.}  
  \label{tab:securebits_values}
\end{table}

A ciascuno dei flag di tab.~\ref{tab:securebits_values} è inoltre abbinato un
corrispondente flag di blocco, identificato da una costante omonima con
l'estensione \texttt{\_LOCKED}, la cui attivazione è irreversibile ed ha
l'effetto di rendere permanente l'impostazione corrente del corrispondente
flag ordinario; in sostanza con \constd{SECURE\_KEEP\_CAPS\_LOCKED} si rende
non più modificabile \const{SECURE\_KEEP\_CAPS}, ed analogamente avviene con
\constd{SECURE\_NO\_SETUID\_FIXUP\_LOCKED} per
\const{SECURE\_NO\_SETUID\_FIXUP} e con \constd{SECURE\_NOROOT\_LOCKED} per
\const{SECURE\_NOROOT}.

Per l'impostazione di questi flag sono state predisposte due specifiche
operazioni di \func{prctl} (vedi sez.~\ref{sec:process_prctl}),
\const{PR\_GET\_SECUREBITS}, che consente di ottenerne il valore, e
\const{PR\_SET\_SECUREBITS}, che consente di modificarne il valore; per
quest'ultima sono comunque necessari i privilegi di amministratore ed in
particolare la capacità \const{CAP\_SETPCAP}. Prima dell'introduzione dei
\textit{securebits} era comunque possibile ottenere lo stesso effetto di
\const{SECURE\_KEEP\_CAPS} attraverso l'uso di un'altra operazione di
\func{prctl}, \const{PR\_SET\_KEEPCAPS}.

\itindend{securebits}

Oltre alla gestione dei \textit{securebits} la nuova versione delle
\textit{file capabilities} prevede l'uso di \func{prctl} anche per la gestione
del \textit{capabilities bounding set}, attraverso altre due operazioni
dedicate, \const{PR\_CAPBSET\_READ} per controllarne il valore e
\const{PR\_CAPBSET\_DROP} per modificarlo; quest'ultima di nuovo è una
operazione privilegiata che richiede la capacità \const{CAP\_SETPCAP} e che,
come indica chiaramente il nome, permette solo la rimozione di una
\textit{capability} dall'insieme; per i dettagli sull'uso di tutte queste
operazioni si rimanda alla rilettura di sez.~\ref{sec:process_prctl}.

\itindend{file~capabilities}


% NOTE per dati relativi al process capability bounding set, vedi:
% http://git.kernel.org/git/?p=linux/kernel/git/torvalds/linux-2.6.git;a=commit;h=3b7391de67da515c91f48aa371de77cb6cc5c07e

% NOTE riferimenti ai vari cambiamenti vedi:
% http://lwn.net/Articles/280279/  
% http://lwn.net/Articles/256519/
% http://lwn.net/Articles/211883/


Un elenco delle \textit{capabilities} disponibili su Linux, con una breve
descrizione ed il nome delle costanti che le identificano, è riportato in
tab.~\ref{tab:proc_capabilities};\footnote{l'elenco presentato questa tabella,
  ripreso dalla pagina di manuale (accessibile con \texttt{man capabilities})
  e dalle definizioni in \texttt{include/linux/capabilities.h}, è aggiornato
  al kernel 3.2.} la tabella è divisa in due parti, la prima riporta le
\textit{capabilities} previste anche nella bozza dello standard POSIX1.e, la
seconda quelle specifiche di Linux.  Come si può notare dalla tabella alcune
\textit{capabilities} attengono a singole funzionalità e sono molto
specializzate, mentre altre hanno un campo di applicazione molto vasto, che è
opportuno dettagliare maggiormente.

\begin{table}[!hbtp]
  \centering
  \footnotesize
  \begin{tabular}{|l|p{10cm}|}
    \hline
    \textbf{Capacità}&\textbf{Descrizione}\\
    \hline
    \hline
%
% POSIX-draft defined capabilities.
%
    \constd{CAP\_AUDIT\_CONTROL}& Abilitare e disabilitare il
                              controllo dell'auditing (dal kernel 2.6.11).\\ 
    \constd{CAP\_AUDIT\_WRITE}&Scrivere dati nel giornale di
                              auditing del kernel (dal kernel 2.6.11).\\ 
    % TODO verificare questa roba dell'auditing
    \constd{CAP\_BLOCK\_SUSPEND}&Utilizzare funzionalità che possono bloccare 
                              la sospensione del sistema (dal kernel 3.5).\\ 
    \constd{CAP\_CHOWN}     & Cambiare proprietario e gruppo
                              proprietario di un file (vedi
                              sez.~\ref{sec:file_ownership_management}).\\
    \constd{CAP\_DAC\_OVERRIDE}& Evitare il controllo dei
                               permessi di lettura, scrittura ed esecuzione dei
                               file, (vedi sez.~\ref{sec:file_access_control}).\\ 
    \constd{CAP\_DAC\_READ\_SEARCH}& Evitare il controllo dei
                              permessi di lettura ed esecuzione per
                              le directory (vedi
                              sez.~\ref{sec:file_access_control}).\\
    \const{CAP\_FOWNER}     & Evitare il controllo della proprietà di un file
                              per tutte le operazioni privilegiate non coperte
                              dalle precedenti \const{CAP\_DAC\_OVERRIDE} e
                              \const{CAP\_DAC\_READ\_SEARCH}.\\
    \constd{CAP\_FSETID}    & Evitare la cancellazione automatica dei bit
                              \acr{suid} e \acr{sgid} quando un file
                              per i quali sono impostati viene modificato da
                              un processo senza questa capacità e la capacità
                              di impostare il bit \acr{sgid} su un file anche
                              quando questo è relativo ad un gruppo cui non si
                              appartiene (vedi
                              sez.~\ref{sec:file_perm_management}).\\ 
    \constd{CAP\_KILL}      & Mandare segnali a qualunque
                              processo (vedi sez.~\ref{sec:sig_kill_raise}).\\
    \constd{CAP\_SETFCAP}   & Impostare le \textit{capabilities} di un file
                              (dal kernel 2.6.24).\\ 
    \constd{CAP\_SETGID}    & Manipolare i group ID dei
                              processi, sia il principale che i supplementari,
                              (vedi sez.~\ref{sec:proc_setgroups}) che quelli
                              trasmessi tramite i socket \textit{unix domain}
                              (vedi sez.~\ref{sec:unix_socket}).\\
    \constd{CAP\_SETUID}    & Manipolare gli user ID del
                              processo (vedi sez.~\ref{sec:proc_setuid}) e di
                              trasmettere un user ID arbitrario nel passaggio
                              delle credenziali coi socket \textit{unix
                                domain} (vedi sez.~\ref{sec:unix_socket}).\\ 
%
% Linux specific capabilities
%
\hline
    \constd{CAP\_IPC\_LOCK} & Effettuare il \textit{memory locking} con le
                              funzioni \func{mlock}, \func{mlockall},
                              \func{shmctl}, \func{mmap} (vedi
                              sez.~\ref{sec:proc_mem_lock} e 
                              sez.~\ref{sec:file_memory_map}). \\ 
% TODO verificare l'interazione con SHM_HUGETLB
    \constd{CAP\_IPC\_OWNER}& Evitare il controllo dei permessi
                              per le operazioni sugli oggetti di
                              intercomunicazione fra processi (vedi
                              sez.~\ref{sec:ipc_sysv}).\\  
    \constd{CAP\_LEASE}     & Creare dei \textit{file lease} (vedi
                              sez.~\ref{sec:file_asyncronous_lease})
                              pur non essendo proprietari del file (dal kernel
                              2.4).\\ 
    \constd{CAP\_LINUX\_IMMUTABLE}& Impostare sui file gli attributi 
                             \textit{immutable} e \textit{append-only} (vedi
                             sez.~\ref{sec:file_perm_overview}) se
                             supportati.\\
    \constd{CAP\_MAC\_ADMIN}& Amministrare il \textit{Mandatory
                               Access Control} di \textit{Smack} (dal kernel
                              2.6.25).\\
    \constd{CAP\_MAC\_OVERRIDE}& Evitare il \textit{Mandatory
                               Access Control} di \textit{Smack} (dal kernel
                              2.6.25).\\   
    \constd{CAP\_MKNOD}     & Creare file di dispositivo con \func{mknod} (vedi
                              sez.~\ref{sec:file_mknod}) (dal kernel 2.4).\\ 
    \const{CAP\_NET\_ADMIN} & Eseguire alcune operazioni
                              privilegiate sulla rete.\\
    \constd{CAP\_NET\_BIND\_SERVICE}& Porsi in ascolto su porte riservate (vedi 
                              sez.~\ref{sec:TCP_func_bind}).\\ 
    \constd{CAP\_NET\_BROADCAST}& Consentire l'uso di socket in
                              \textit{broadcast} e \textit{multicast}.\\ 
    \constd{CAP\_NET\_RAW}  & Usare socket \texttt{RAW} e \texttt{PACKET}
                              (vedi sez.~\ref{sec:sock_type}).\\ 
    \const{CAP\_SETPCAP}    & Effettuare modifiche privilegiate alle
                              \textit{capabilities}.\\   
    \const{CAP\_SYS\_ADMIN} & Eseguire una serie di compiti amministrativi.\\
    \constd{CAP\_SYS\_BOOT} & Eseguire un riavvio del sistema (vedi
                              sez.~\ref{sec:sys_reboot}).\\ 
    \constd{CAP\_SYS\_CHROOT}& Eseguire la funzione \func{chroot} (vedi 
                              sez.~\ref{sec:file_chroot}).\\
    \constd{CAP\_SYS\_MODULE}& Caricare e rimuovere moduli del kernel.\\ 
    \const{CAP\_SYS\_NICE}  & Modificare le varie priorità dei processi (vedi 
                              sez.~\ref{sec:proc_priority}).\\
    \constd{CAP\_SYS\_PACCT}& Usare le funzioni di \textit{accounting} dei 
                              processi (vedi
                              sez.~\ref{sec:sys_bsd_accounting}).\\  
    \constd{CAP\_SYS\_PTRACE}& La capacità di tracciare qualunque processo con
                              \func{ptrace} (vedi 
                              sez.~\ref{sec:process_ptrace}).\\
    \constd{CAP\_SYS\_RAWIO}& Operare sulle porte di I/O con \func{ioperm} e
                               \func{iopl} (vedi
                              sez.~\ref{sec:process_io_port}).\\
    \const{CAP\_SYS\_RESOURCE}& Superare le varie limitazioni sulle risorse.\\ 
    \constd{CAP\_SYS\_TIME} & Modificare il tempo di sistema (vedi 
                              sez.~\ref{sec:sys_time}).\\ 
    \constd{CAP\_SYS\_TTY\_CONFIG}&Simulare un \textit{hangup} della console,
                              con la funzione \func{vhangup}.\\
    \constd{CAP\_SYSLOG}    & Gestire il buffer dei messaggi
                              del kernel, (vedi sez.~\ref{sec:sess_daemon}),
                              introdotta dal kernel 2.6.38 come capacità
                              separata da \const{CAP\_SYS\_ADMIN}.\\
    \constd{CAP\_WAKE\_ALARM}&Usare i timer di tipo
                              \const{CLOCK\_BOOTTIME\_ALARM} e
                              \const{CLOCK\_REALTIME\_ALARM}, vedi
                              sez.~\ref{sec:sig_timer_adv} (dal kernel 3.0).\\  
    \hline
  \end{tabular}
  \caption{Le costanti che identificano le \textit{capabilities} presenti nel
    kernel.}
\label{tab:proc_capabilities}
\end{table}

\constbeg{CAP\_SETPCAP}

Prima di dettagliare il significato della capacità più generiche, conviene
però dedicare un discorso a parte a \const{CAP\_SETPCAP}, il cui significato è
stato completamente cambiato con l'introduzione delle \textit{file
  capabilities} nel kernel 2.6.24. In precedenza questa capacità era quella
che permetteva al processo che la possedeva di impostare o rimuovere le
\textit{capabilities} presenti nel suo \textit{permitted set} su un qualunque
altro processo. In realtà questo non è mai stato l'uso inteso nelle bozze
dallo standard POSIX, ed inoltre, come si è già accennato, dato che questa
capacità è sempre stata assente (a meno di specifiche ricompilazioni del
kernel) nel \textit{capabilities bounding set} usato di default, essa non è
neanche mai stata realmente disponibile.

Con l'introduzione \textit{file capabilities} e il cambiamento del significato
del \textit{capabilities bounding set} la possibilità di modificare le
capacità di altri processi è stata completamente rimossa, e
\const{CAP\_SETPCAP} ha acquisito quello che avrebbe dovuto essere il suo
significato originario, e cioè la capacità del processo di poter inserire nel
suo \textit{inheritable set} qualunque capacità presente nel \textit{bounding
  set}. Oltre a questo la disponibilità di \const{CAP\_SETPCAP} consente ad un
processo di eliminare una capacità dal proprio \textit{bounding set} (con la
conseguente impossibilità successiva di eseguire programmi con quella
capacità), o di impostare i \textit{securebits} delle \textit{capabilities}.

\constend{CAP\_SETPCAP}
\constbeg{CAP\_FOWNER}

La prima fra le capacità ``\textsl{ampie}'' che occorre dettagliare
maggiormente è \const{CAP\_FOWNER}, che rimuove le restrizioni poste ad un
processo che non ha la proprietà di un file in un vasto campo di
operazioni;\footnote{vale a dire la richiesta che l'\ids{UID} effettivo del
  processo (o meglio l'\ids{UID} di filesystem, vedi
  sez.~\ref{sec:proc_setuid}) coincida con quello del proprietario.}  queste
comprendono i cambiamenti dei permessi e dei tempi del file (vedi
sez.~\ref{sec:file_perm_management} e sez.~\ref{sec:file_file_times}), le
impostazioni degli attributi dei file e delle ACL (vedi
sez.~\ref{sec:file_xattr} e \ref{sec:file_ACL}), poter ignorare lo
\textit{sticky bit} nella cancellazione dei file (vedi
sez.~\ref{sec:file_special_perm}), la possibilità di impostare il flag di
\const{O\_NOATIME} con \func{open} e \func{fcntl} (vedi
sez.~\ref{sec:file_open_close} e sez.~\ref{sec:file_fcntl_ioctl}) senza
restrizioni.

\constend{CAP\_FOWNER}
\constbeg{CAP\_NET\_ADMIN}

Una seconda capacità che copre diverse operazioni, in questo caso riguardanti
la rete, è \const{CAP\_NET\_ADMIN}, che consente di impostare le opzioni
privilegiate dei socket (vedi sez.~\ref{sec:sock_generic_options}), abilitare
il \textit{multicasting} (vedi sez.\ref{sec:sock_ipv4_options}), eseguire la
configurazione delle interfacce di rete (vedi
sez.~\ref{sec:sock_ioctl_netdevice}) ed impostare la tabella di instradamento.

\constend{CAP\_NET\_ADMIN}
\constbeg{CAP\_SYS\_ADMIN}

Una terza \textit{capability} con vasto campo di applicazione è
\const{CAP\_SYS\_ADMIN}, che copre una serie di operazioni amministrative,
come impostare le quote disco (vedi sez.\ref{sec:disk_quota}), attivare e
disattivare la \textit{swap}, montare, rimontare e smontare filesystem (vedi
sez.~\ref{sec:filesystem_mounting}), effettuare operazioni di controllo su
qualunque oggetto dell'IPC di SysV (vedi sez.~\ref{sec:ipc_sysv}), operare
sugli attributi estesi dei file di classe \texttt{security} o \texttt{trusted}
(vedi sez.~\ref{sec:file_xattr}), specificare un \ids{UID} arbitrario nella
trasmissione delle credenziali dei socket (vedi
sez.~\ref{sec:socket_credential_xxx}), assegnare classi privilegiate
(\const{IOPRIO\_CLASS\_RT} e prima del kernel 2.6.25 anche
\const{IOPRIO\_CLASS\_IDLE}) per lo scheduling dell'I/O (vedi
sez.~\ref{sec:io_priority}), superare il limite di sistema sul numero massimo
di file aperti,\footnote{quello indicato da \sysctlfiled{fs/file-max}.}
effettuare operazioni privilegiate sulle chiavi mantenute dal kernel (vedi
sez.~\ref{sec:keyctl_management}), usare la funzione \func{lookup\_dcookie},
usare \const{CLONE\_NEWNS} con \func{unshare} e \func{clone}, (vedi
sez.~\ref{sec:process_clone}).

\constend{CAP\_SYS\_ADMIN}
\constbeg{CAP\_SYS\_NICE}

Originariamente \const{CAP\_SYS\_NICE} riguardava soltanto la capacità di
aumentare le priorità di esecuzione dei processi, come la diminuzione del
valore di \textit{nice} (vedi sez.~\ref{sec:proc_sched_stand}), l'uso delle
priorità \textit{real-time} (vedi sez.~\ref{sec:proc_real_time}), o
l'impostazione delle affinità di processore (vedi
sez.~\ref{sec:proc_sched_multiprocess}); ma con l'introduzione di priorità
anche riguardo le operazioni di accesso al disco, e, nel caso di sistemi NUMA,
alla memoria, essa viene a coprire anche la possibilità di assegnare priorità
arbitrarie nell'accesso a disco (vedi sez.~\ref{sec:io_priority}) e nelle
politiche di allocazione delle pagine di memoria ai nodi di un sistema NUMA.

\constend{CAP\_SYS\_NICE}
\constbeg{CAP\_SYS\_RESOURCE}

Infine la \textit{capability} \const{CAP\_SYS\_RESOURCE} attiene alla
possibilità di superare i limiti imposti sulle risorse di sistema, come usare
lo spazio disco riservato all'amministratore sui filesystem che lo supportano,
usare la funzione \func{ioctl} per controllare il \textit{journaling} sul
filesystem \acr{ext3}, non subire le quote disco, aumentare i limiti sulle
risorse di un processo (vedi sez.~\ref{sec:sys_resource_limit}) e quelle sul
numero di processi, ed i limiti sulle dimensioni dei messaggi delle code del
SysV IPC (vedi sez.~\ref{sec:ipc_sysv_mq}).

\constend{CAP\_SYS\_RESOURCE}

Per la gestione delle \textit{capabilities} il kernel mette a disposizione due
funzioni che permettono rispettivamente di leggere ed impostare i valori dei
tre insiemi illustrati in precedenza. Queste due funzioni di sistema sono
\funcd{capget} e \funcd{capset} e costituiscono l'interfaccia di gestione
basso livello; i loro rispettivi prototipi sono:

\begin{funcproto}{
\fhead{sys/capability.h}
\fdecl{int capget(cap\_user\_header\_t hdrp, cap\_user\_data\_t datap)}
\fdesc{Legge le \textit{capabilities}.} 
\fdecl{int capset(cap\_user\_header\_t hdrp, const cap\_user\_data\_t datap)} 
\fdesc{Imposta le \textit{capabilities}.} 
}

{Le funzioni ritornano $0$ in caso di successo e $-1$ per un errore, nel qual
  caso \var{errno} assumerà uno dei valori: 
  \begin{errlist}
  \item[\errcode{EFAULT}] si è indicato un puntatore sbagliato o nullo
    per \param{hdrp} o \param{datap} (quest'ultimo può essere nullo solo se si
    usa \func{capget} per ottenere la versione delle \textit{capabilities}
    usata dal kernel).
  \item[\errcode{EINVAL}] si è specificato un valore non valido per uno dei
    campi di \param{hdrp}, in particolare una versione non valida della
    versione delle \textit{capabilities}.
  \item[\errcode{EPERM}] si è tentato di aggiungere una capacità nell'insieme
    delle \textit{capabilities} permesse, o di impostare una capacità non
    presente nell'insieme di quelle permesse negli insieme delle effettive o
    ereditate, o si è cercato di impostare una \textit{capability} di un altro
    processo senza avare \const{CAP\_SETPCAP}.
  \item[\errcode{ESRCH}] si è fatto riferimento ad un processo inesistente.
  \end{errlist}
}
\end{funcproto}

Queste due funzioni prendono come argomenti due tipi di dati dedicati,
definiti come puntatori a due strutture specifiche di Linux, illustrate in
fig.~\ref{fig:cap_kernel_struct}.  Per un certo periodo di tempo era anche
indicato che per poterle utilizzare fosse necessario che la macro
\macro{\_POSIX\_SOURCE} risultasse non definita (ed era richiesto di inserire
una istruzione \texttt{\#undef \_POSIX\_SOURCE} prima di includere
\headfiled{sys/capability.h}) requisito che non risulta più
presente.\footnote{e non è chiaro neanche quanto sia mai stato davvero
  necessario.}

\begin{figure}[!htb]
  \footnotesize
  \centering
  \begin{minipage}[c]{0.8\textwidth}
    \includestruct{listati/cap_user_header_t.h}
  \end{minipage}
  \normalsize 
  \caption{Definizione delle strutture a cui fanno riferimento i puntatori
    \structd{cap\_user\_header\_t} e \structd{cap\_user\_data\_t} usati per
    l'interfaccia di gestione di basso livello delle \textit{capabilities}.}
  \label{fig:cap_kernel_struct}
\end{figure}

Si tenga presente che le strutture di fig.~\ref{fig:cap_kernel_struct}, come i
prototipi delle due funzioni \func{capget} e \func{capset}, sono soggette ad
essere modificate con il cambiamento del kernel (in particolare i tipi di dati
delle strutture) ed anche se finora l'interfaccia è risultata stabile, non c'è
nessuna assicurazione che questa venga mantenuta,\footnote{viene però
  garantito che le vecchie funzioni continuino a funzionare.} Pertanto se si
vogliono scrivere programmi portabili che possano essere eseguiti senza
modifiche o adeguamenti su qualunque versione del kernel è opportuno
utilizzare le interfacce di alto livello che vedremo più avanti.

La struttura a cui deve puntare l'argomento \param{hdrp} serve ad indicare,
tramite il campo \var{pid}, il \ids{PID} del processo del quale si vogliono
leggere o modificare le \textit{capabilities}. Con \func{capset} questo, se si
usano le \textit{file capabilities}, può essere solo 0 o il \ids{PID} del
processo chiamante, che sono equivalenti. Non tratteremo, essendo comunque di
uso irrilevante, il caso in cui, in mancanza di tale supporto, la funzione può
essere usata per modificare le \textit{capabilities} di altri processi, per il
quale si rimanda, se interessati, alla lettura della pagina di manuale.

Il campo \var{version} deve essere impostato al valore della versione delle
stesse usata dal kernel (quello indicato da una delle costanti
\texttt{\_LINUX\_CAPABILITY\_VERSION\_n} di fig.~\ref{fig:cap_kernel_struct})
altrimenti le funzioni ritorneranno con un errore di \errcode{EINVAL},
restituendo nel campo stesso il valore corretto della versione in uso. La
versione due è comunque deprecata e non deve essere usata, ed il kernel
stamperà un avviso se lo si fa.

I valori delle \textit{capabilities} devono essere passati come maschere
binarie;\footnote{e si tenga presente che i valori di
  tab.~\ref{tab:proc_capabilities} non possono essere combinati direttamente,
  indicando il numero progressivo del bit associato alla relativa capacità.}
con l'introduzione delle \textit{capabilities} a 64 bit inoltre il
puntatore \param{datap} non può essere più considerato come relativo ad una
singola struttura, ma ad un vettore di due strutture.\footnote{è questo cambio
  di significato che ha portato a deprecare la versione 2, che con
  \func{capget} poteva portare ad un buffer overflow per vecchie applicazioni
  che continuavano a considerare \param{datap} come puntatore ad una singola
  struttura.}

Dato che le precedenti funzioni, oltre ad essere specifiche di Linux, non
garantiscono la stabilità nell'interfaccia, è sempre opportuno effettuare la
gestione delle \textit{capabilities} utilizzando le funzioni di libreria a
questo dedicate. Queste funzioni, che seguono quanto previsto nelle bozze
dello standard POSIX.1e, non fanno parte della \acr{glibc} e sono fornite in
una libreria a parte,\footnote{la libreria è \texttt{libcap2}, nel caso di
  Debian può essere installata con il pacchetto omonimo.} pertanto se un
programma le utilizza si dovrà indicare esplicitamente al compilatore l'uso
della suddetta libreria attraverso l'opzione \texttt{-lcap}.

\itindbeg{capability~state}

Le funzioni dell'interfaccia alle \textit{capabilities} definite nelle bozze
dello standard POSIX.1e prevedono l'uso di un tipo di dato opaco,
\typed{cap\_t}, come puntatore ai dati mantenuti nel cosiddetto
\textit{capability state},\footnote{si tratta in sostanza di un puntatore ad
  una struttura interna utilizzata dalle librerie, i cui campi non devono mai
  essere acceduti direttamente.} in sono memorizzati tutti i dati delle
\textit{capabilities}.

In questo modo è possibile mascherare i dettagli della gestione di basso
livello, che potranno essere modificati senza dover cambiare le funzioni
dell'interfaccia, che fanno riferimento soltanto ad oggetti di questo tipo.
L'interfaccia pertanto non soltanto fornisce le funzioni per modificare e
leggere le \textit{capabilities}, ma anche quelle per gestire i dati
attraverso i \textit{capability state}, che presentano notevoli affinità,
essendo parte di bozze dello stesso standard, con quelle già viste per le ACL.

La prima funzione dell'interfaccia è quella che permette di inizializzare un
\textit{capability state}, allocando al contempo la memoria necessaria per i
relativi dati. La funzione è \funcd{cap\_init} ed il suo prototipo è:

\begin{funcproto}{
\fhead{sys/capability.h}
\fdecl{cap\_t cap\_init(void)}
\fdesc{Crea ed inizializza un \textit{capability state}.} 
}

{La funzione ritorna un \textit{capability state} in caso di successo e
  \val{NULL} per un errore, nel qual caso \var{errno} potrà assumere solo il
  valore \errval{ENOMEM}.  }
\end{funcproto}

La funzione restituisce il puntatore \type{cap\_t} ad uno stato inizializzato
con tutte le \textit{capabilities} azzerate. In caso di errore (cioè quando
non c'è memoria sufficiente ad allocare i dati) viene restituito \val{NULL}
ed \var{errno} viene impostata a \errval{ENOMEM}.  

La memoria necessaria a mantenere i dati viene automaticamente allocata da
\func{cap\_init}, ma dovrà essere disallocata esplicitamente quando non è più
necessaria utilizzando, per questo l'interfaccia fornisce una apposita
funzione, \funcd{cap\_free}, il cui prototipo è:

\begin{funcproto}{
\fhead{sys/capability.h}
\fdecl{int cap\_free(void *obj\_d)}
\fdesc{Disalloca la memoria allocata per i dati delle \textit{capabilities}..} 
}

{La funzione ritorna $0$ in caso di successo e $-1$ per un errore, nel qual
  caso \var{errno} potrà assumere solo il valore \errval{EINVAL}.
}
\end{funcproto}


La funzione permette di liberare la memoria allocata dalle altre funzioni
della libreria sia per un \textit{capability state}, nel qual caso l'argomento
sarà un dato di tipo \type{cap\_t}, che per una descrizione testuale dello
stesso,\footnote{cioè quanto ottenuto tramite la funzione
  \func{cap\_to\_text}.} nel qual caso l'argomento sarà un dato di tipo
\texttt{char *}. Per questo motivo l'argomento \param{obj\_d} è dichiarato
come \texttt{void *}, per evitare la necessità di eseguire un \textit{cast},
ma dovrà comunque corrispondere ad un puntatore ottenuto tramite le altre
funzioni della libreria, altrimenti la funzione fallirà con un errore di
\errval{EINVAL}.

Infine si può creare una copia di un \textit{capability state} ottenuto in
precedenza tramite la funzione \funcd{cap\_dup}, il cui prototipo è:

\begin{funcproto}{
\fhead{sys/capability.h}
\fdecl{cap\_t cap\_dup(cap\_t cap\_p)}
\fdesc{Duplica un \textit{capability state} restituendone una copia.} 
}

{La funzione ritorna un \textit{capability state} in caso di successo e
  \val{NULL} per un errore, nel qual caso \var{errno} assumerà i valori
  \errval{ENOMEM} o \errval{EINVAL} nel loro significato generico.}
\end{funcproto}


La funzione crea una copia del \textit{capability state} posto all'indirizzo
\param{cap\_p} che si è passato come argomento, restituendo il puntatore alla
copia, che conterrà gli stessi valori delle \textit{capabilities} presenti
nell'originale. La memoria necessaria viene allocata automaticamente dalla
funzione. Una volta effettuata la copia i due \textit{capability state}
potranno essere modificati in maniera completamente indipendente, ed alla fine
delle operazioni si dovrà disallocare anche la copia, oltre all'originale.

Una seconda classe di funzioni di servizio previste dall'interfaccia sono
quelle per la gestione dei dati contenuti all'interno di un \textit{capability
  state}; la prima di queste è \funcd{cap\_clear}, il cui prototipo è:

\begin{funcproto}{
\fhead{sys/capability.h}
\fdecl{int cap\_clear(cap\_t cap\_p)}
\fdesc{Inizializza un \textit{capability state} cancellando tutte le
  \textit{capabilities}.}
}

{La funzione ritorna $0$ in caso di successo e $-1$ per un errore, nel qual
  caso \var{errno} potrà assumere solo il valore \errval{EINVAL}.
}
\end{funcproto}

La funzione si limita ad azzerare tutte le \textit{capabilities} presenti nel
\textit{capability state} all'indirizzo \param{cap\_p} passato come argomento,
restituendo uno stato \textsl{vuoto}, analogo a quello che si ottiene nella
creazione con \func{cap\_init}.

Una variante di \func{cap\_clear} è \funcd{cap\_clear\_flag} che cancella da
un \textit{capability state} tutte le \textit{capabilities} di un certo
insieme fra quelli elencati a pag.~\pageref{sec:capabilities_set}, il suo
prototipo è:

\begin{funcproto}{
\fhead{sys/capability.h}
\fdecl{int cap\_clear\_flag(cap\_t cap\_p, cap\_flag\_t flag)} 
\fdesc{Cancella delle \textit{capabilities} da un \textit{capability state}.} 
}

{La funzione ritorna $0$ in caso di successo e $-1$ per un errore, nel qual
  caso \var{errno}  potrà assumere solo il valore \errval{EINVAL}.
}
\end{funcproto}

La funzione richiede che si indichi quale degli insiemi si intente cancellare
da \param{cap\_p} con l'argomento \param{flag}. Questo deve essere specificato
con una variabile di tipo \type{cap\_flag\_t} che può assumere
esclusivamente\footnote{si tratta in effetti di un tipo enumerato, come si può
  verificare dalla sua definizione che si trova in
  \headfile{sys/capability.h}.} uno dei valori illustrati in
tab.~\ref{tab:cap_set_identifier}.

\begin{table}[htb]
  \centering
  \footnotesize
  \begin{tabular}[c]{|l|l|}
    \hline
    \textbf{Valore} & \textbf{Significato} \\
    \hline
    \hline
    \constd{CAP\_EFFECTIVE}  & Capacità dell'insieme \textsl{effettivo}.\\
    \constd{CAP\_PERMITTED}  & Capacità dell'insieme \textsl{permesso}.\\ 
    \constd{CAP\_INHERITABLE}& Capacità dell'insieme \textsl{ereditabile}.\\
    \hline
  \end{tabular}
  \caption{Valori possibili per il tipo di dato \typed{cap\_flag\_t} che
    identifica gli insiemi delle \textit{capabilities}.}
  \label{tab:cap_set_identifier}
\end{table}

Si possono inoltre confrontare in maniera diretta due diversi
\textit{capability state} con la funzione \funcd{cap\_compare}; il suo
prototipo è:

\begin{funcproto}{
\fhead{sys/capability.h}
\fdecl{int cap\_compare(cap\_t cap\_a, cap\_t cap\_b)}
\fdesc{Confronta due \textit{capability state}.} 
}

{La funzione ritorna $0$ se i \textit{capability state} sono identici
    ed un valore positivo se differiscono, non sono previsti errori.}
\end{funcproto}


La funzione esegue un confronto fra i due \textit{capability state} passati
come argomenti e ritorna in un valore intero il risultato, questo è nullo se
sono identici o positivo se vi sono delle differenze. Il valore di ritorno
della funzione consente inoltre di per ottenere ulteriori informazioni su
quali sono gli insiemi di \textit{capabilities} che risultano differenti.  Per
questo si può infatti usare la apposita macro \macro{CAP\_DIFFERS}:

{\centering
\vspace{3pt}
\begin{funcbox}{
\fhead{sys/capability.h}
\fdecl{int \macrod{CAP\_DIFFERS}(value, flag)}
\fdesc{Controlla lo stato di eventuali differenze delle \textit{capabilities}
  nell'insieme \texttt{flag}.}
}
\end{funcbox}
}

La macro richiede che si passi nell'argomento \texttt{value} il risultato
della funzione \func{cap\_compare} e in \texttt{flag} l'indicazione (coi
valori di tab.~\ref{tab:cap_set_identifier}) dell'insieme che si intende
controllare; restituirà un valore diverso da zero se le differenze rilevate da
\func{cap\_compare} sono presenti nell'insieme indicato.

Per la gestione dei singoli valori delle \textit{capabilities} presenti in un
\textit{capability state} l'interfaccia prevede due funzioni specifiche,
\funcd{cap\_get\_flag} e \funcd{cap\_set\_flag}, che permettono
rispettivamente di leggere o impostare il valore di una capacità all'interno
in uno dei tre insiemi già citati; i rispettivi prototipi sono:

\begin{funcproto}{
\fhead{sys/capability.h}
\fdecl{int cap\_get\_flag(cap\_t cap\_p, cap\_value\_t cap, cap\_flag\_t 
flag,\\
\phantom{int cap\_get\_flag(}cap\_flag\_value\_t *value\_p)}
\fdesc{Legge il valore di una \textit{capability}.}
\fdecl{int cap\_set\_flag(cap\_t cap\_p, cap\_flag\_t flag, int ncap,
  cap\_value\_t *caps, \\
\phantom{int cap\_set\_flag(}cap\_flag\_value\_t value)} 
\fdesc{Imposta il valore di una \textit{capability}.} 
}

{Le funzioni ritornano $0$ in caso di successo e $-1$ per un errore, nel qual
  caso \var{errno} potrà assumere solo il valore \errval{EINVAL}.  
}
\end{funcproto}

In entrambe le funzioni l'argomento \param{cap\_p} indica il puntatore al
\textit{capability state} su cui operare, mentre l'argomento \param{flag}
indica su quale dei tre insiemi si intende operare, sempre con i valori di
tab.~\ref{tab:cap_set_identifier}.  La capacità che si intende controllare o
impostare invece deve essere specificata attraverso una variabile di tipo
\typed{cap\_value\_t}, che può prendere come valore uno qualunque di quelli
riportati in tab.~\ref{tab:proc_capabilities}, in questo caso però non è
possibile combinare diversi valori in una maschera binaria, una variabile di
tipo \type{cap\_value\_t} può indicare una sola capacità.\footnote{in
  \headfile{sys/capability.h} il tipo \type{cap\_value\_t} è definito come
  \ctyp{int}, ma i valori validi sono soltanto quelli di
  tab.~\ref{tab:proc_capabilities}.}

Infine lo stato di una capacità è descritto ad una variabile di tipo
\type{cap\_flag\_value\_t}, che a sua volta può assumere soltanto
uno\footnote{anche questo è un tipo enumerato.} dei valori di
tab.~\ref{tab:cap_value_type}.

\begin{table}[htb]
  \centering
  \footnotesize
  \begin{tabular}[c]{|l|l|}
    \hline
    \textbf{Valore} & \textbf{Significato} \\
    \hline
    \hline
    \constd{CAP\_CLEAR}& La capacità non è impostata.\\ 
    \constd{CAP\_SET}  & La capacità è impostata.\\
    \hline
  \end{tabular}
  \caption{Valori possibili per il tipo di dato \typed{cap\_flag\_value\_t} che
    indica lo stato di una capacità.}
  \label{tab:cap_value_type}
\end{table}

La funzione \func{cap\_get\_flag} legge lo stato della capacità indicata
dall'argomento \param{cap} all'interno dell'insieme indicato dall'argomento
\param{flag} e lo restituisce come \textit{value result argument} nella
variabile puntata dall'argomento \param{value\_p}. Questa deve essere di tipo
\type{cap\_flag\_value\_t} ed assumerà uno dei valori di
tab.~\ref{tab:cap_value_type}. La funzione consente pertanto di leggere solo
lo stato di una capacità alla volta.

La funzione \func{cap\_set\_flag} può invece impostare in una sola chiamata
più \textit{capabilities}, anche se solo all'interno dello stesso insieme ed
allo stesso valore. Per questo motivo essa prende un vettore di valori di tipo
\type{cap\_value\_t} nell'argomento \param{caps}, la cui dimensione viene
specificata dall'argomento \param{ncap}. Il tipo di impostazione da eseguire
(cancellazione o attivazione) per le capacità elencate in \param{caps} viene
indicato dall'argomento \param{value} sempre con i valori di
tab.~\ref{tab:cap_value_type}.

Per semplificare la gestione delle \textit{capabilities} l'interfaccia prevede
che sia possibile utilizzare anche una rappresentazione testuale del contenuto
di un \textit{capability state} e fornisce le opportune funzioni di
gestione;\footnote{entrambe erano previste dalla bozza dello standard
  POSIX.1e.} la prima di queste, che consente di ottenere la rappresentazione
testuale, è \funcd{cap\_to\_text}, il cui prototipo è:

\begin{funcproto}{
\fhead{sys/capability.h}
\fdecl{char *cap\_to\_text(cap\_t caps, ssize\_t *length\_p)}
\fdesc{Genera una visualizzazione testuale delle \textit{capabilities}.} 
}

{La funzione ritorna un puntatore alla stringa con la descrizione delle
  \textit{capabilities} in caso di successo e \val{NULL} per un errore, nel
  qual caso \var{errno} assumerà i valori \errval{EINVAL} o \errval{ENOMEM}
  nel loro significato generico.}
\end{funcproto}

La funzione ritorna l'indirizzo di una stringa contente la descrizione
testuale del contenuto del \textit{capability state} \param{caps} passato come
argomento, e, qualora l'argomento \param{length\_p} sia diverso da \val{NULL},
restituisce come \textit{value result argument} nella variabile intera da
questo puntata la lunghezza della stringa. La stringa restituita viene
allocata automaticamente dalla funzione e pertanto dovrà essere liberata con
\func{cap\_free}.

La rappresentazione testuale, che viene usata anche dai programmi di gestione a
riga di comando, prevede che lo stato venga rappresentato con una stringa di
testo composta da una serie di proposizioni separate da spazi, ciascuna delle
quali specifica una operazione da eseguire per creare lo stato finale. Nella
rappresentazione si fa sempre conto di partire da uno stato in cui tutti gli
insiemi sono vuoti e si provvede a impostarne i contenuti.

Ciascuna proposizione è nella forma di un elenco di capacità, espresso con i
nomi di tab.~\ref{tab:proc_capabilities} separati da virgole, seguito da un
operatore, e dall'indicazione degli insiemi a cui l'operazione si applica. I
nomi delle capacità possono essere scritti sia maiuscoli che minuscoli, viene
inoltre riconosciuto il nome speciale \texttt{all} che è equivalente a
scrivere la lista completa. Gli insiemi sono identificati dalle tre lettere
iniziali: ``\texttt{p}'' per il \textit{permitted}, ``\texttt{i}'' per
l'\textit{inheritable} ed ``\texttt{e}'' per l'\textit{effective} che devono
essere sempre minuscole, e se ne può indicare più di uno.

Gli operatori possibili sono solo tre: ``\texttt{+}'' che aggiunge le capacità
elencate agli insiemi indicati, ``\texttt{-}'' che le toglie e ``\texttt{=}''
che le assegna esattamente. I primi due richiedono che sia sempre indicato sia
un elenco di capacità che gli insiemi a cui esse devono applicarsi, e
rispettivamente attiveranno o disattiveranno le capacità elencate nell'insieme
o negli insiemi specificati, ignorando tutto il resto. I due operatori possono
anche essere combinati nella stessa proposizione, per aggiungere e togliere le
capacità dell'elenco da insiemi diversi.

L'assegnazione si applica invece su tutti gli insiemi allo stesso tempo,
pertanto l'uso di ``\texttt{=}'' è equivalente alla cancellazione preventiva
di tutte le capacità ed alla impostazione di quelle elencate negli insiemi
specificati, questo significa che in genere lo si usa una sola volta
all'inizio della stringa. In tal caso l'elenco delle capacità può non essere
indicato e viene assunto che si stia facendo riferimento a tutte quante senza
doverlo scrivere esplicitamente.

Come esempi avremo allora che un processo non privilegiato di un utente, che
non ha nessuna capacità attiva, avrà una rappresentazione nella forma
``\texttt{=}'' che corrisponde al fatto che nessuna capacità viene assegnata a
nessun insieme (vale la cancellazione preventiva), mentre un processo con
privilegi di amministratore avrà una rappresentazione nella forma
``\texttt{=ep}'' in cui tutte le capacità vengono assegnate agli insiemi
\textit{permitted} ed \textit{effective} (e l'\textit{inheritable} è ignorato
in quanto per le regole viste a pag.~\ref{sec:capability-uid-transition} le
capacità verranno comunque attivate attraverso una \func{exec}). Infine, come
esempio meno banale dei precedenti, otterremo per \texttt{init} una
rappresentazione nella forma ``\texttt{=ep cap\_setpcap-e}'' dato che come
accennato tradizionalmente \const{CAP\_SETPCAP} è sempre stata rimossa da
detto processo.

Viceversa per ottenere un \textit{capability state} dalla sua rappresentazione
testuale si può usare la funzione \funcd{cap\_from\_text}, il cui prototipo è:

\begin{funcproto}{
\fhead{sys/capability.h}
\fdecl{cap\_t cap\_from\_text(const char *string)}
\fdesc{Crea un \textit{capability state} dalla sua rappresentazione testuale.} 
}

{La funzione ritorna un \textit{capability state} in caso di successo e
  \val{NULL} per un errore, nel qual caso \var{errno} assumerà i valori
  \errval{EINVAL} o \errval{ENOMEM} nel loro significato generico.}
\end{funcproto}


La funzione restituisce il puntatore ad un \textit{capability state}
inizializzato con i valori indicati nella stringa \param{string} che ne
contiene la rappresentazione testuale. La memoria per il \textit{capability
  state} viene allocata automaticamente dalla funzione e dovrà essere liberata
con \func{cap\_free}.

Alle due funzioni citate se ne aggiungono altre due che consentono di
convertire i valori delle costanti di tab.~\ref{tab:proc_capabilities} nelle
stringhe usate nelle rispettive rappresentazioni e viceversa. Le due funzioni,
\funcd{cap\_to\_name} e \funcd{cap\_from\_name}, sono estensioni specifiche di
Linux ed i rispettivi prototipi sono:

\begin{funcproto}{
\fhead{sys/capability.h}
\fdecl{char *cap\_to\_name(cap\_value\_t cap)}
\fdesc{Converte il valore numerico di una \textit{capabilities} alla sua
  rappresentazione testuale.} 
\fdecl{int cap\_from\_name(const char *name, cap\_value\_t *cap\_p)}

\fdesc{Converte la rappresentazione testuale di una \textit{capabilities} al
  suo valore numerico.} 
}

{La funzione \func{cap\_to\_name} ritorna un puntatore ad una stringa in caso
  di successo e \val{NULL} per un errore, mentre \func{cap\_to\_name} ritorna
  $0$ in caso di successo e $-1$ per un errore, per entrambe in caso di errore
  \var{errno} assumerà i valori \errval{EINVAL} o \errval{ENOMEM} nel loro
  significato generico.  
}
\end{funcproto}

La prima funzione restituisce la stringa (allocata automaticamente e che dovrà
essere liberata con \func{cap\_free}) che corrisponde al valore della
capacità \param{cap}, mentre la seconda restituisce nella variabile puntata
da \param{cap\_p}, come \textit{value result argument}, il valore della
capacità rappresentata dalla stringa \param{name}.

Fin quei abbiamo trattato solo le funzioni di servizio relative alla
manipolazione dei \textit{capability state} come strutture di dati;
l'interfaccia di gestione prevede però anche le funzioni per trattare le
\textit{capabilities} presenti nei processi. La prima di queste funzioni è
\funcd{cap\_get\_proc} che consente la lettura delle \textit{capabilities} del
processo corrente, il suo prototipo è:

\begin{funcproto}{
\fhead{sys/capability.h}
\fdecl{cap\_t cap\_get\_proc(void)}
\fdesc{Legge le \textit{capabilities} del processo corrente.} 
}

{La funzione ritorna un \textit{capability state} in caso di successo e
  \val{NULL} per un errore, nel qual caso \var{errno} assumerà i valori
  \errval{EINVAL}, \errval{EPERM} o \errval{ENOMEM} nel loro significato
  generico.}
\end{funcproto}

La funzione legge il valore delle \textit{capabilities} associate al processo
da cui viene invocata, restituendo il risultato tramite il puntatore ad un
\textit{capability state} contenente tutti i dati che provvede ad allocare
autonomamente e che di nuovo occorrerà liberare con \func{cap\_free} quando
non sarà più utilizzato.

Se invece si vogliono leggere le \textit{capabilities} di un processo
specifico occorre usare la funzione \funcd{cap\_get\_pid}, il cui
prototipo\footnote{su alcune pagine di manuale la funzione è descritta con un
  prototipo sbagliato, che prevede un valore di ritorno di tipo \type{cap\_t},
  ma il valore di ritorno è intero, come si può verificare anche dalla
  dichiarazione della stessa in \headfile{sys/capability.h}.} è:

\begin{funcproto}{
\fhead{sys/capability.h}
\fdecl{cap\_t cap\_get\_pid(pid\_t pid)}
\fdesc{Legge le \textit{capabilities} di un processo.} 
}

{La funzione ritorna un \textit{capability state} in caso di successo e
  \val{NULL} per un errore, nel qual caso \var{errno} assumerà i valori
  \errval{ESRCH} o \errval{ENOMEM} nel loro significato generico.  }
\end{funcproto}

La funzione legge il valore delle \textit{capabilities} del processo indicato
con l'argomento \param{pid}, e restituisce il risultato tramite il puntatore
ad un \textit{capability state} contenente tutti i dati che provvede ad
allocare autonomamente e che al solito deve essere disallocato con
\func{cap\_free}. Qualora il processo indicato non esista si avrà un errore di
\errval{ESRCH}. Gli stessi valori possono essere letti direttamente nel
filesystem \textit{proc}, nei file \texttt{/proc/<pid>/status}; ad esempio per
\texttt{init} si otterrà qualcosa del tipo:
\begin{Console}
piccardi@hain:~/gapil$ \textbf{cat /proc/1/status}
...
CapInh: 0000000000000000
CapPrm: 00000000fffffeff
CapEff: 00000000fffffeff  
...
\end{Console}
%$

\itindend{capability~state}

Infine per impostare le \textit{capabilities} del processo corrente (nella
bozza dello standard POSIX.1e non esiste una funzione che permetta di cambiare
le \textit{capabilities} di un altro processo) si deve usare la funzione
\funcd{cap\_set\_proc}, il cui prototipo è:

\begin{funcproto}{
\fhead{sys/capability.h}
\fdecl{int cap\_set\_proc(cap\_t cap\_p)}
\fdesc{Imposta le \textit{capabilities} del processo corrente.} 
}

{La funzione ritorna $0$ in caso di successo e $-1$ per un errore, nel qual
  caso \var{errno} assumerà i valori:
  \begin{errlist}
  \item[\errcode{EPERM}] si è cercato di attivare una capacità non permessa.
  \end{errlist} ed inoltre \errval{EINVAL} nel suo significato generico.}
\end{funcproto}

La funzione modifica le \textit{capabilities} del processo corrente secondo
quanto specificato con l'argomento \param{cap\_p}, posto che questo sia
possibile nei termini spiegati in precedenza (non sarà ad esempio possibile
impostare capacità non presenti nell'insieme di quelle permesse). 

In caso di successo i nuovi valori saranno effettivi al ritorno della
funzione, in caso di fallimento invece lo stato delle capacità resterà
invariato. Si tenga presente che \textsl{tutte} le capacità specificate
tramite \param{cap\_p} devono essere permesse; se anche una sola non lo è la
funzione fallirà, e per quanto appena detto, lo stato delle
\textit{capabilities} non verrà modificato (neanche per le parti eventualmente
permesse).

Oltre a queste funzioni su Linux sono presenti due ulteriori funzioni,
\funcm{capgetp} e \funcm{capsetp}, che svolgono un compito analogo. Queste
funzioni risalgono alla implementazione iniziale delle \textit{capabilities}
ed in particolare \funcm{capsetp} consentirebbe anche, come possibile in quel
caso, di cambiare le capacità di un altro processo. Le due funzioni oggi sono
deprecate e pertanto eviteremo di trattarle, per chi fosse interessato si
rimanda alla lettura della loro pagina di manuale.

Come esempio di utilizzo di queste funzioni nei sorgenti allegati alla guida
si è distribuito il programma \texttt{getcap.c}, che consente di leggere le
\textit{capabilities} del processo corrente\footnote{vale a dire di sé stesso,
  quando lo si lancia, il che può sembrare inutile, ma serve a mostrarci quali
  sono le \textit{capabilities} standard che ottiene un processo lanciato
  dalla riga di comando.} o tramite l'opzione \texttt{-p}, quelle di un
processo qualunque il cui \ids{PID} viene passato come parametro dell'opzione.

\begin{figure}[!htbp]
  \footnotesize \centering
  \begin{minipage}[c]{\codesamplewidth}
    \includecodesample{listati/getcap.c}
  \end{minipage} 
  \normalsize
  \caption{Corpo principale del programma \texttt{getcap.c}.}
  \label{fig:proc_getcap}
\end{figure}

La sezione principale del programma è riportata in fig.~\ref{fig:proc_getcap},
e si basa su una condizione sulla variabile \var{pid} che se si è usato
l'opzione \texttt{-p} è impostata (nella sezione di gestione delle opzioni,
che si è tralasciata) al valore del \ids{PID} del processo di cui si vuole
leggere le \textit{capabilities} e nulla altrimenti. Nel primo caso
(\texttt{\small 1-6}) si utilizza (\texttt{\small 2}) \func{cap\_get\_proc}
per ottenere lo stato delle capacità del processo, nel secondo (\texttt{\small
  7-13}) si usa invece \func{cap\_get\_pid} (\texttt{\small 8}) per leggere
il valore delle capacità del processo indicato.

Il passo successivo è utilizzare (\texttt{\small 15}) \func{cap\_to\_text} per
tradurre in una stringa lo stato, e poi (\texttt{\small 16}) stamparlo; infine
(\texttt{\small 18-19}) si libera la memoria allocata dalle precedenti
funzioni con \func{cap\_free} per poi ritornare dal ciclo principale della
funzione.

\itindend{capabilities}

% TODO vedi http://lwn.net/Articles/198557/ e 
% http://www.madore.org/~david/linux/newcaps/




\subsection{La gestione del \textit{Secure Computing}.}
\label{sec:procadv_seccomp}

\itindbeg{secure~computing~mode}

Il \textit{secure computing mode} è un meccanismo ideato per fornire un
supporto per l'esecuzione di codice esterno non fidato e non verificabile a
scopo di calcolo. L'idea era quella di disporre di una modalità di esecuzione
dei programmi che permettesse di vendere la capacità di calcolo della propria
macchina ad un qualche servizio di calcolo distribuito, senza comprometterne
la sicurezza eseguendo codice non sotto il proprio controllo.

La prima versione del meccanismo è stata introdotta con il kernel
2.6.23,\footnote{e disponibile solo avendo abilitato il supporto nel kernel
  con l'opzione di configurazione \texttt{CONFIG\_SECCOMP}.} è molto semplice,
il \textit{secure computing mode} viene attivato con \func{prctl} usando
l'opzione \const{PR\_SET\_SECCOMP}, ed indicando \const{SECCOMP\_MODE\_STRICT}
come valore per \param{arg2} (all'epoca unico valore possibile).  Una volta
abilitato in questa modalità (in seguito denominata \textit{strict mode}) il
processo o il \textit{thread} chiamante potrà utilizzare soltanto un insieme
estremamente limitato di \textit{system call}: \func{read}, \func{write},
\func{\_exit} e \funcm{sigreturn}; l'esecuzione di qualsiasi altra
\textit{system call} comporta l'emissione di un \signal{SIGKILL} e conseguente
terminazione immediata del processo.

Si tenga presente che in questo caso, con versioni recenti della \acr{glibc}
(il comportamento è stato introdotto con la 2.3), diventa impossibile usare
anche \func{\_exit} in \textit{strict mode}, in quanto questa funzione viene
intercettata ed al suo posto viene chiamata \func{exit\_group} (vedi
sez.~\ref{sec:pthread_management}) che non è consentita e comporta un
\signal{SIGKILL}.

Si tenga presente che, non essendo \func{execve} fra le funzioni permesse, per
poter eseguire un programma terzo essendo in \textit{strict mode} questo dovrà
essere fornito in una forma di codice interpretabile fornito attraverso un
socket o una \textit{pipe}, creati prima di lanciare il processo che eseguirà
il codice non fidato. 


% TODO a partire dal kernel 3.5 è stato introdotto la possibilità di usare un
% terzo argomento se il secondo è SECCOMP_MODE_FILTER, vedi
% Documentation/prctl/seccomp_filter.txt 
% vedi anche http://lwn.net/Articles/600250/

% TODO documentare PR_SET_SECCOMP introdotto a partire dal kernel 3.5. Vedi:
% * Documentation/prctl/seccomp_filter.txt
% * http://lwn.net/Articles/475043/

% TODO a partire dal kernel 3.17 è stata introdotta la nuova syscall seccomp,
% vedi http://lwn.net/Articles/600250/ e http://lwn.net/Articles/603321/


\itindend{secure~computing~mode}

\subsection{Altre funzionalità di sicurezza.}
\label{sec:procadv_security_misc}

Oltre alle funzionalità specifiche esaminate nelle sezioni precedenti, il
kernel supporta una varietà di ulteriori impostazioni di sicurezza,
accessibili nelle maniere più varie, che abbiamo raccolto in questa sezione.

Una serie di modalità di sicurezza sono attivabili a richiesta attraverso
alcune opzioni di controllo attivabili via \func{sysctl} o il filesystem
\texttt{/proc}, un elenco delle stesse e dei loro effetti è il seguente:

\begin{basedescript}{\desclabelwidth{1cm}\desclabelstyle{\nextlinelabel}}
\item[\sysctlrelfiled{fs}{protected\_hardlinks}] Un valore nullo, il default,
  mantiene il comportamento standard che non pone restrizioni alla creazione
  di \textit{hard link}. Se il valore viene posto ad 1 vengono invece attivate
  una serie di restrizioni protettive, denominate
  \itindex{protected~hardlinks} \textit{protected hardlinks}, che se non
  soddisfatte causano il fallimento di \func{link} con un errore di
  \errval{EPERM}. Perché questo non avvenga almeno una delle seguenti
  condizioni deve essere soddisfatta:
  \begin{itemize*}
  \item il chiamante deve avere privilegi amministrativi (la
    \textit{capability} \const{CAP\_FOWNER}). In caso di utilizzo
    dell'\textit{user namespace} oltre a possedere \const{CAP\_FOWNER} è
    necessario che l'\ids{UID} del proprietario del file sia mappato nel
    \textit{namespace}.
  \item il \textit{filesystem} \ids{UID} del chiamante (normalmente
    equivalente all'\ids{UID} effettivo) deve corrispondere a quello del
    proprietario del file a cui si vuole effettuare il collegamento.
  \item devono essere soddisfatte tutte le seguenti condizioni:
    \begin{itemize*}
    \item il file è un file ordinario
    \item il file non ha il \acr{suid} bit attivo
    \item il file non ha lo \acr{sgid} bit attivo ed il permesso di esecuzione
      per il gruppo
    \item il chiamante ha i permessi di lettura e scrittura sul file
    \end{itemize*}
  \end{itemize*}

  In sostanza in questo caso un utente potrà creare un collegamento diretto ad
  un altro file solo se ne è il proprietario o se questo è un file ordinario
  senza permessi speciali ed a cui ha accesso in lettura e scrittura.

  Questa funzionalità fornisce una protezione generica che non inficia l'uso
  ordinario di \func{link}, ma rende impraticabili una serie di possibili
  abusi della stessa; oltre ad impedire l'uso di un \textit{hard link} come
  variante in un attacco di \textit{symlink race} (eludendo i
  \textit{protected symlinks} di cui al punto successivo), evita anche che si
  possa lasciare un riferimento ad un eventuale programma \acr{suid}
  vulnerabile, creando un collegamento diretto allo stesso.


\item[\sysctlrelfiled{fs}{protected\_symlinks}] Un valore nullo, il default,
  mantiene il comportamento standard che non pone restrizioni nel seguire i
  link simbolici. Se il valore viene posto ad 1 vengono attivate delle
  restrizioni protettive, denominate \itindex{protected~symlinks}
  \textit{protected symlinks}. Quando vengono attivate una qualunque funzione
  che esegua la risoluzione di un \textit{pathname} contenente un link
  simbolico non conforme alle restrizioni fallirà con un errore di
  \errval{EACCESS}. Per evitare l'errore deve essere soddisfatta una delle
  seguenti condizioni:
  \begin{itemize*}
  \item il link non è in una directory con permessi analoghi a \file{/tmp}
    (scrivibile a tutti e con lo \textit{sticky bit} attivo);
  \item il link è in una directory con permessi analoghi a \file{/tmp} ma è
    soddisfatta una delle condizioni seguenti: 
    \begin{itemize*}
    \item il link simbolico appartiene al chiamante: il controllo viene fatto
      usando il \textit{filesystem} \ids{UID} (che normalmente corrisponde
      all'\ids{UID} effettivo).
    \item il link simbolico ha lo stesso proprietario della directory.
    \end{itemize*}
  \end{itemize*}

  Questa funzionalità consente di rendere impraticabili alcuni attacchi in cui
  si approfitta di una differenza di tempo fra il controllo e l'uso di un
  file, ed in particolare quella classe di attacchi viene usualmente chiamati
  \textit{symlink attack},\footnote{si tratta di un sottoinsieme di quella
    classe di attacchi chiamata genericamente \textit{TOCTTOU}, acronimo
    appunto di \textit{Time of check to time of use}.} di cui abbiamo parlato
  in sez.~\ref{sec:file_temp_file}.

  Un possibile esempio di questo tipo di attacco è quello contro un programma
  che viene eseguito per conto di un utente privilegiato (ad esempio un
  programma con il \acr{suid} o lo \acr{sgid} bit attivi) che prima controlla
  l'esistenza di un file e se non esiste lo crea. Se questa procedura, che è
  tipica della creazione di file temporanei sotto \file{/tmp}, non viene
  eseguita in maniera corretta,\footnote{ad esempio con le modalità che
    abbiamo trattato in sez.~\ref{sec:file_temp_file}, che per quanto note da
    tempo continuano ad essere ignorate.} un attaccante ha una finestra di
  tempo in cui può creare prima del programma un \textit{link simbolico} ad un
  file di sua scelta, compresi file di dispositivo o file a cui non avrebbe
  accesso, facendolo poi utilizzare al programma.

  Attivando la funzionalità si rende impossibile seguire un link simbolico in
  una directory temporanea come \texttt{/tmp}, a meno che questo non sia di
  proprietà del chiamante, o che questo non appartenga al proprietario della
  directory. Questo impedisce che i link simbolici creati da un attaccante
  possano essere seguiti da un programma privilegiato (perché apparterranno
  all'attaccante) mentre quelli creati dall'amministratore (che i genere è il
  proprietario di \texttt{/tmp}) saranno seguiti comunque.

\end{basedescript}


% TODO: trattare pure protected_regular e protected_fifos introdotti con il
% 4.19 (vedi https://lwn.net/Articles/763106/)



% TODO: trattare keyctl (man 2 keyctl)



% TODO trattare le funzioni di protezione della memoria pkey_alloc, pkey_free,
% pkey_mprotect, introdotte con il kernel 4.8, vedi
% http://lwn.net/Articles/689395/ e Documentation/x86/protection-keys.txt 

\section{Funzioni di gestione e controllo}
\label{sec:proc_manage_control}

In questa sezione prenderemo in esame alcune specifiche \textit{system call}
dedicate al controllo dei processi sia per quanto riguarda l'impostazione di
caratteristiche specialistiche, che per quanto riguarda l'analisi ed il
controllo della loro esecuzione.

\subsection{La funzione \func{prctl}}
\label{sec:process_prctl}

Benché la gestione ordinaria dei processi possa essere effettuata attraverso
le funzioni che abbiamo già esaminato nei capitoli \ref{cha:process_interface}
e \ref{cha:process_handling}, esistono una serie di proprietà e
caratteristiche specifiche dei processi per la cui gestione è stata
predisposta una apposita \textit{system call} che fornisce una interfaccia
generica per tutte le operazioni specialistiche. La funzione di sistema è
\funcd{prctl} ed il suo prototipo è:\footnote{la funzione non è standardizzata
  ed è specifica di Linux, anche se ne esiste una analoga in IRIX; è stata
  introdotta con il kernel 2.1.57.}

\begin{funcproto}{ 
\fhead{sys/prctl.h}
\fdecl{int prctl(int option, unsigned long arg2, unsigned long arg3, unsigned
  long arg4, \\
\phantom{int prctl(}unsigned long arg5)}
\fdesc{Esegue una operazione speciale sul processo corrente.} 
}

{La funzione ritorna $0$ o un valore positivo dipendente dall'operazione in
  caso di successo e $-1$ per un errore, nel qual caso \var{errno} assumerà
  valori diversi a seconda del tipo di operazione richiesta, sono possibili:
  \errval{EACCESS}, \errval{EBADF}, \errval{EBUSY}, \errval{EFAULT},
  \errval{EINVAL}, \errval{ENXIO}, \errval{EOPNOTSUPP} o \errval{EPERM}.}
\end{funcproto}

La funzione ritorna in caso di successo un valore nullo o positivo, e $-1$ in
caso di errore. Il significato degli argomenti della funzione successivi al
primo, il valore di ritorno in caso di successo, il tipo di errore restituito
in \var{errno} dipendono dall'operazione eseguita, indicata tramite il primo
argomento, \param{option}. Questo è un valore intero che identifica
l'operazione, e deve essere specificato con l'uso di una delle costanti
predefinite del seguente elenco.\footnote{l'elenco potrebbe non risultare
  aggiornato, in quanto nuove operazioni vengono aggiunte nello sviluppo del
  kernel.} Tratteremo esplicitamente per ciascuna di esse il significato del
il valore di ritorno in caso di successo, ma solo quando non corrisponde
all'ordinario valore nullo (dato per implicito).

%TODO: trattare PR_CAP_AMBIENT, dal 4.3
%TODO: trattare PR_CAP_FP_*, dal 4.0, solo per MIPS
%TODO: trattare PR_MPX_*_MANAGEMENT, dal 3.19
%TODO: trattare PR_*NO_NEW_PRIVS, dal 3.5

\begin{basedescript}{\desclabelwidth{1.5cm}\desclabelstyle{\nextlinelabel}}
\item[\constd{PR\_CAPBSET\_READ}] Controlla la disponibilità di una delle
  \textit{capability} (vedi sez.~\ref{sec:proc_capabilities}). La funzione
  ritorna 1 se la capacità specificata nell'argomento \param{arg2} (con una
  delle costanti di tab.~\ref{tab:proc_capabilities}) è presente nel
  \textit{capabilities bounding set} del processo e zero altrimenti,
  se \param{arg2} non è un valore valido si avrà un errore di \errval{EINVAL}.
  Introdotta a partire dal kernel 2.6.25.

\item[\constd{PR\_CAPBSET\_DROP}] Rimuove permanentemente una delle
  \textit{capabilities} (vedi sez.~\ref{sec:proc_capabilities}) dal processo e
  da tutti i suoi discendenti. La funzione cancella la capacità specificata
  nell'argomento \param{arg2} con una delle costanti di
  tab.~\ref{tab:proc_capabilities} dal \textit{capabilities bounding set} del
  processo. L'operazione richiede i privilegi di amministratore (la capacità
  \const{CAP\_SETPCAP}), altrimenti la chiamata fallirà con un errore di
  \errcode{EPERM}; se il valore di \param{arg2} non è valido o se il supporto
  per le \textit{file capabilities} non è stato compilato nel kernel la
  chiamata fallirà con un errore di \errval{EINVAL}. Introdotta a partire dal
  kernel 2.6.25.

\item[\constd{PR\_SET\_DUMPABLE}] Imposta il flag che determina se la
  terminazione di un processo a causa di un segnale per il quale è prevista la
  generazione di un file di \textit{core dump} (vedi
  sez.~\ref{sec:sig_standard}) lo genera effettivamente. In genere questo flag
  viene attivato automaticamente, ma per evitare problemi di sicurezza (la
  generazione di un file da parte di processi privilegiati può essere usata
  per sovrascriverne altri) viene cancellato quando si mette in esecuzione un
  programma con i bit \acr{suid} e \acr{sgid} attivi (vedi
  sez.~\ref{sec:file_special_perm}) o con l'uso delle funzioni per la modifica
  degli \ids{UID} dei processi (vedi sez.~\ref{sec:proc_setuid}).

  L'operazione è stata introdotta a partire dal kernel 2.3.20, fino al kernel
  2.6.12 e per i kernel successivi al 2.6.17 era possibile usare solo un
  valore 0 (espresso anche come \constd{SUID\_DUMP\_DISABLE}) di \param{arg2}
  per disattivare il flag ed un valore 1 (espresso anche come
  \constd{SUID\_DUMP\_USER}) per attivarlo. Nei kernel dal 2.6.13 al 2.6.17 è
  stato supportato anche il valore 2, che causava la generazione di un
  \textit{core dump} leggibile solo dall'amministratore, ma questa
  funzionalità è stata rimossa per motivi di sicurezza, in quanto consentiva
  ad un utente normale di creare un file di \textit{core dump} appartenente
  all'amministratore in directory dove l'utente avrebbe avuto permessi di
  accesso. Specificando un valore diverso da 0 o 1 si ottiene un errore di
  \errval{EINVAL}.

\item[\constd{PR\_GET\_DUMPABLE}] Ottiene come valore di ritorno della funzione
  lo stato corrente del flag che controlla la effettiva generazione dei
  \textit{core dump}. Introdotta a partire dal kernel 2.3.20.

\item[\constd{PR\_SET\_ENDIAN}] Imposta la \textit{endianness} del processo
  chiamante secondo il valore fornito in \param{arg2}. I valori possibili sono
  sono: \constd{PR\_ENDIAN\_BIG} (\textit{big endian}),
  \constd{PR\_ENDIAN\_LITTLE} (\textit{little endian}), e
  \constd{PR\_ENDIAN\_PPC\_LITTLE} (lo pseudo \textit{little endian} del
  PowerPC). Introdotta a partire dal kernel 2.6.18, solo per architettura
  PowerPC.

\item[\constd{PR\_GET\_ENDIAN}] Ottiene il valore della \textit{endianness} del
  processo chiamante, salvato sulla variabile puntata da \param{arg2} che deve
  essere passata come di tipo ``\ctyp{int *}''. Introdotta a partire dal
  kernel 2.6.18, solo su PowerPC.

\item[\constd{PR\_SET\_FPEMU}] Imposta i bit di controllo per l'emulazione
  della virgola mobile su architettura ia64, secondo il valore
  di \param{arg2}, si deve passare \constd{PR\_FPEMU\_NOPRINT} per emulare in
  maniera trasparente l'accesso alle operazioni in virgola mobile, o
  \constd{PR\_FPEMU\_SIGFPE} per non emularle ed inviare il segnale
  \signal{SIGFPE} (vedi sez.~\ref{sec:sig_prog_error}). Introdotta a partire
  dal kernel 2.4.18, solo su architettura ia64.

\item[\constd{PR\_GET\_FPEMU}] Ottiene il valore dei flag di controllo
  dell'emulazione della virgola mobile, salvato all'indirizzo puntato
  da \param{arg2}, che deve essere di tipo ``\ctyp{int *}''. Introdotta a
  partire dal kernel 2.4.18, solo su architettura ia64.

\item[\constd{PR\_SET\_FPEXC}] Imposta la modalità delle eccezioni in virgola
  mobile (\textit{floating-point exception mode}) al valore di \param{arg2}.
  I valori possibili sono: 
  \begin{itemize*}
  \item \constd{PR\_FP\_EXC\_SW\_ENABLE} per usare FPEXC per le eccezioni,
  \item \constd{PR\_FP\_EXC\_DIV} per la divisione per zero in virgola mobile,
  \item \constd{PR\_FP\_EXC\_OVF} per gli overflow,
  \item \constd{PR\_FP\_EXC\_UND} per gli underflow,
  \item \constd{PR\_FP\_EXC\_RES} per risultati non esatti,
  \item \constd{PR\_FP\_EXC\_INV} per operazioni invalide,
  \item \constd{PR\_FP\_EXC\_DISABLED} per disabilitare le eccezioni,
  \item \constd{PR\_FP\_EXC\_NONRECOV} per usare la modalità di eccezione
    asincrona non recuperabile,
  \item \constd{PR\_FP\_EXC\_ASYNC} per usare la modalità di eccezione
    asincrona recuperabile,
  \item \constd{PR\_FP\_EXC\_PRECISE} per la modalità precisa di
    eccezione.\footnote{trattasi di gestione specialistica della gestione
      delle eccezioni dei calcoli in virgola mobile che, i cui dettagli al
      momento vanno al di là dello scopo di questo testo.}
  \end{itemize*}
Introdotta a partire dal kernel 2.4.21, solo su PowerPC.

\item[\constd{PR\_GET\_FPEXC}] Ottiene il valore della modalità delle eccezioni
  delle operazioni in virgola mobile, salvata all'indirizzo
  puntato \param{arg2}, che deve essere di tipo ``\ctyp{int *}''.  Introdotta
  a partire dal kernel 2.4.21, solo su PowerPC.

\item[\constd{PR\_SET\_KEEPCAPS}] Consente di controllare quali
  \textit{capabilities} vengono cancellate quando si esegue un cambiamento di
  \ids{UID} del processo (per i dettagli si veda
  sez.~\ref{sec:proc_capabilities}, in particolare quanto illustrato a
  pag.~\pageref{sec:capability-uid-transition}). Un valore nullo (il default)
  per \param{arg2} comporta che vengano cancellate, il valore 1 che vengano
  mantenute, questo valore viene sempre cancellato attraverso una \func{exec}.
  L'uso di questo flag è stato sostituito, a partire dal kernel 2.6.26, dal
  flag \const{SECURE\_KEEP\_CAPS} dei \textit{securebits} (vedi
  sez.~\ref{sec:proc_capabilities} e l'uso di \const{PR\_SET\_SECUREBITS} più
  avanti) e si è impostato con essi \const{SECURE\_KEEP\_CAPS\_LOCKED} si
  otterrà un errore di \errval{EPERM}.  Introdotta a partire dal kernel
  2.2.18.

\item[\constd{PR\_GET\_KEEPCAPS}] Ottiene come valore di ritorno della funzione
  il valore del flag di controllo delle \textit{capabilities} impostato con
  \const{PR\_SET\_KEEPCAPS}. Introdotta a partire dal kernel 2.2.18.

\item[\constd{PR\_SET\_NAME}] Imposta il nome del processo chiamante alla
  stringa puntata da \param{arg2}, che deve essere di tipo ``\ctyp{char *}''. Il
  nome può essere lungo al massimo 16 caratteri, e la stringa deve essere
  terminata da NUL se più corta.  Introdotta a partire dal kernel 2.6.9.

\item[\constd{PR\_GET\_NAME}] Ottiene il nome del processo chiamante nella
  stringa puntata da \param{arg2}, che deve essere di tipo ``\ctyp{char *}'';
  si devono allocare per questo almeno 16 byte, e il nome sarà terminato da
  NUL se più corto. Introdotta a partire dal kernel 2.6.9.

\item[\constd{PR\_SET\_PDEATHSIG}] Consente di richiedere l'emissione di un
  segnale, che sarà ricevuto dal processo chiamante, in occorrenza della
  terminazione del proprio processo padre; in sostanza consente di invertire
  il ruolo di \signal{SIGCHLD}. Il valore di \param{arg2} deve indicare il
  numero del segnale, o 0 per disabilitare l'emissione. Il valore viene
  automaticamente cancellato per un processo figlio creato con \func{fork}.
  Introdotta a partire dal kernel 2.1.57.

\item[\constd{PR\_GET\_PDEATHSIG}] Ottiene il valore dell'eventuale segnale
  emesso alla terminazione del padre, salvato all'indirizzo
  puntato \param{arg2}, che deve essere di tipo ``\ctyp{int *}''. Introdotta a
  partire dal kernel 2.3.15.

\item[\constd{PR\_SET\_PTRACER}] Imposta un \ids{PID} per il ``\textit{tracer
    process}'' usando \param{arg2}. Una impostazione successiva sovrascrive la
  precedente, ed un valore nullo cancella la disponibilità di un
  ``\textit{tracer process}''. Questa è una funzionalità fornita da
  \textit{``Yama''}, uno specifico \textit{Linux Security Modules}, e serve a
  consentire al processo indicato, quando le restrizioni introdotte da questo
  modulo sono attive, di usare \func{ptrace} (vedi
  sez.~\ref{sec:process_ptrace}) sul processo chiamante, anche se quello
  indicato non ne è un progenitore. Il valore \constd{PR\_SET\_PTRACER\_ANY}
  consente a tutti i processi l'uso di \func{ptrace}. L'uso si \textit{Yama}
  attiene alla gestione della sicurezza dei processi, e consente di introdurre
  una restrizione all'uso di \func{ptrace}, che è spesso sorgente di
  compromissioni. Si tratta di un uso specialistico che va al di là dello
  scopo di queste dispense, per i dettagli si consulti la documentazione su
  \textit{Yama} nei sorgenti del kernel. Introdotta a partire dal kernel 3.4.

\item[\constd{PR\_SET\_SECCOMP}] Attiva il \textit{secure computing mode} per
  il processo corrente. Introdotta a partire dal kernel 2.6.23 la funzionalità
  è stata ulteriormente estesa con il kernel 3.5, salvo poi diventare un
  sottoinsieme della \textit{system call} \func{seccomp} a partire dal kernel
  3.17. Prevede che si indichi per \param{arg2} il valore
  \const{SECCOMP\_MODE\_STRICT} (unico possibile fino al kernel 2.6.23) per
  selezionare il cosiddetto \textit{strict mode} o, dal kernel 3.5,
  \const{SECCOMP\_MODE\_FILTER} per usare il \textit{filter mode}. Tratteremo
  questa opzione nei dettagli più avanti, in sez.~\ref{sec:procadv_seccomp},
  quando affronteremo l'argomento del \textit{Secure Computing}.

\item[\constd{PR\_GET\_SECCOMP}] Ottiene come valore di ritorno della funzione
  lo stato corrente del \textit{secure computing mode}. Fino al kernel 3.5,
  quando era possibile solo lo \textit{strict mode}, la funzione era
  totalmente inutile in quanto l'unico valore ottenibile era 0 in assenza di
  \textit{secure computing}, dato che la chiamata di questa funzione in
  \textit{strict mode} avrebbe comportato l'emissione di \signal{SIGKILL} per
  il chiamante. La funzione però, a partire dal kernel 2.6.23, era stata
  comunque definita per eventuali estensioni future, ed infatti con
  l'introduzione del \textit{filter mode} con il kernel 3.5, se essa viene
  inclusa nelle funzioni consentite restituisce il valore 2 quando il
  \textit{secure computing mode} è attivo (se non inclusa si avrà di nuovo un
  \signal{SIGKILL}).

\item[\constd{PR\_SET\_SECUREBITS}] Imposta i \textit{securebits} per il
  processo chiamante al valore indicato da \param{arg2}; per i dettagli sul
  significato dei \textit{securebits} si veda
  sez.~\ref{sec:proc_capabilities}, ed in particolare i valori di
  tab.~\ref{tab:securebits_values} e la relativa trattazione. L'operazione
  richiede i privilegi di amministratore (la capacità \const{CAP\_SETPCAP}),
  altrimenti la chiamata fallirà con un errore di \errval{EPERM}. Introdotta a
  partire dal kernel 2.6.26.

\item[\constd{PR\_GET\_SECUREBITS}] Ottiene come valore di ritorno della
  funzione l'impostazione corrente per i \textit{securebits}. Introdotta a
  partire dal kernel 2.6.26.

\item[\constd{PR\_SET\_TIMING}] Imposta il metodo di temporizzazione del
  processo da indicare con il valore di \param{arg2}, attualmente i valori
  possibili sono due, con \constd{PR\_TIMING\_STATISTICAL} si usa il metodo
  statistico tradizionale, con \constd{PR\_TIMING\_TIMESTAMP} il più accurato
  basato su dei \textit{timestamp}, quest'ultimo però non è ancora
  implementato ed il suo uso comporta la restituzione di un errore di
  \errval{EINVAL}. Introdotta a partire dal kernel 2.6.0-test4.

\item[\constd{PR\_GET\_TIMING}] Ottiene come valore di ritorno della funzione
  il metodo di temporizzazione del processo attualmente in uso (uno dei due
  valori citati per \const{PR\_SET\_TIMING}). Introdotta a partire dal kernel
  2.6.0-test4.

\item[\constd{PR\_SET\_TSC}] Imposta il flag che indica se il processo
  chiamante può leggere il registro di processore contenente il contatore dei
  \textit{timestamp} (TSC, o \textit{Time Stamp Counter}) da indicare con il
  valore di \param{arg2}. Si deve specificare \constd{PR\_TSC\_ENABLE} per
  abilitare la lettura o \constd{PR\_TSC\_SIGSEGV} per disabilitarla con la
  generazione di un segnale di \signal{SIGSEGV} (vedi
  sez.~\ref{sec:sig_prog_error}). La lettura viene automaticamente
  disabilitata se si attiva il \textit{secure computing mode} (vedi
  \const{PR\_SET\_SECCOMP} e sez.~\ref{sec:procadv_seccomp}).  Introdotta a
  partire dal kernel 2.6.26, solo su x86.

\item[\constd{PR\_GET\_TSC}] Ottiene il valore del flag che controlla la
  lettura del contattore dei \textit{timestamp}, salvato all'indirizzo
  puntato \param{arg2}, che deve essere di tipo ``\ctyp{int *}''. Introdotta a
  partire dal kernel 2.6.26, solo su x86.
% articoli sul TSC e relativi problemi: http://lwn.net/Articles/209101/,
% http://blog.cr0.org/2009/05/time-stamp-counter-disabling-oddities.html,
% http://en.wikipedia.org/wiki/Time_Stamp_Counter 

\item[\constd{PR\_SET\_UNALIGN}] Imposta la modalità di controllo per l'accesso
  a indirizzi di memoria non allineati, che in varie architetture risultano
  illegali, da indicare con il valore di \param{arg2}. Si deve specificare il
  valore \constd{PR\_UNALIGN\_NOPRINT} per ignorare gli accessi non allineati,
  ed il valore \constd{PR\_UNALIGN\_SIGBUS} per generare un segnale di
  \signal{SIGBUS} (vedi sez.~\ref{sec:sig_prog_error}) in caso di accesso non
  allineato.  Introdotta con diverse versioni su diverse architetture.

\item[\const{PR\_GET\_UNALIGN}] Ottiene il valore della modalità di controllo
  per l'accesso a indirizzi di memoria non allineati, salvato all'indirizzo
  puntato \param{arg2}, che deve essere di tipo \code{(int *)}. Introdotta con
  diverse versioni su diverse architetture.
\item[\const{PR\_MCE\_KILL}] Imposta la politica di gestione degli errori
  dovuti a corruzione della memoria per problemi hardware. Questo tipo di
  errori vengono riportati dall'hardware di controllo della RAM e vengono
  gestiti dal kernel,\footnote{la funzionalità è disponibile solo sulle
    piattaforme più avanzate che hanno il supporto hardware per questo tipo di
    controlli.} ma devono essere opportunamente riportati ai processi che
  usano quella parte di RAM che presenta errori; nel caso specifico questo
  avviene attraverso l'emissione di un segnale di \signal{SIGBUS} (vedi
  sez.~\ref{sec:sig_prog_error}).\footnote{in particolare viene anche
    impostato il valore di \var{si\_code} in \struct{siginfo\_t} a
    \const{BUS\_MCEERR\_AO}; per il significato di tutto questo si faccia
    riferimento alla trattazione di sez.~\ref{sec:sig_sigaction}.}

  Il comportamento di default prevede che per tutti i processi si applichi la
  politica generale di sistema definita nel file
  \sysctlfiled{vm/memory\_failure\_early\_kill}, ma specificando
  per \param{arg2} il valore \constd{PR\_MCE\_KILL\_SET} è possibile impostare
  con il contenuto di \param{arg3} una politica specifica del processo
  chiamante. Si può tornare alla politica di default del sistema utilizzando
  invece per \param{arg2} il valore \constd{PR\_MCE\_KILL\_CLEAR}. In tutti i
  casi, per compatibilità con eventuali estensioni future, tutti i valori
  degli argomenti non utilizzati devono essere esplicitamente posti a zero,
  pena il fallimento della chiamata con un errore di \errval{EINVAL}.
  
  In caso di impostazione di una politica specifica del processo con
  \const{PR\_MCE\_KILL\_SET} i valori di \param{arg3} possono essere soltanto
  due, che corrispondono anche al valore che si trova nell'impostazione
  generale di sistema di \texttt{memory\_failure\_early\_kill}, con
  \constd{PR\_MCE\_KILL\_EARLY} si richiede l'emissione immediata di
  \signal{SIGBUS} non appena viene rilevato un errore, mentre con
  \constd{PR\_MCE\_KILL\_LATE} il segnale verrà inviato solo quando il processo
  tenterà un accesso alla memoria corrotta. Questi due valori corrispondono
  rispettivamente ai valori 1 e 0 di
  \texttt{memory\_failure\_early\_kill}.\footnote{in sostanza nel primo caso
    viene immediatamente inviato il segnale a tutti i processi che hanno la
    memoria corrotta mappata all'interno del loro spazio degli indirizzi, nel
    secondo caso prima la pagina di memoria viene tolta dallo spazio degli
    indirizzi di ciascun processo, mentre il segnale viene inviato solo quei
    processi che tentano di accedervi.} Si può usare per \param{arg3} anche un
  terzo valore, \constd{PR\_MCE\_KILL\_DEFAULT}, che corrisponde a impostare
  per il processo la politica di default.\footnote{si presume la politica di
    default corrente, in modo da non essere influenzati da un eventuale
    successivo cambiamento della stessa.} Introdotta a partire dal kernel
  2.6.32.
\item[\constd{PR\_MCE\_KILL\_GET}] Ottiene come valore di ritorno della
  funzione la politica di gestione degli errori dovuti a corruzione della
  memoria. Tutti gli argomenti non utilizzati (al momento tutti) devono essere
  nulli pena la ricezione di un errore di \errval{EINVAL}. Introdotta a
  partire dal kernel 2.6.32.
\itindbeg{child~reaper}
\item[\constd{PR\_SET\_CHILD\_SUBREAPER}] Se \param{arg2} è diverso da zero
  imposta l'attributo di \textit{child reaper} per il processo, se nullo lo
  cancella. Lo stato di \textit{child reaper} è una funzionalità, introdotta
  con il kernel 3.4, che consente di far svolgere al processo che ha questo
  attributo il ruolo di ``\textsl{genitore adottivo}'' per tutti i processi
  suoi ``\textsl{discendenti}'' che diventano orfani, in questo modo il
  processo potrà ricevere gli stati di terminazione alla loro uscita,
  sostituendo in questo ruolo \cmd{init} (si ricordi quanto illustrato in
  sez.~\ref{sec:proc_termination}). Il meccanismo è stato introdotto ad uso
  dei programmi di gestione dei servizi, per consentire loro di ricevere gli
  stati di terminazione di tutti i processi che lanciano, anche se questi
  eseguono una doppia \func{fork}; nel comportamento ordinario infatti questi
  verrebbero adottati da \cmd{init} ed il programma che li ha lanciati non
  sarebbe più in grado di riceverne lo stato di terminazione. Se un processo
  con lo stato di \textit{child reaper} termina prima dei suoi discendenti,
  svolgerà questo ruolo il più prossimo antenato ad avere lo stato di
  \textit{child reaper}, 
\item[\constd{PR\_GET\_CHILD\_SUBREAPER}] Ottiene l'impostazione relativa allo
  lo stato di \textit{child reaper} del processo chiamante, salvata come
  \textit{value result} all'indirizzo puntato da \param{arg2} (da indicare
  come di tipo \code{int *}). Il valore viene letto come valore logico, se
  diverso da 0 lo stato di \textit{child reaper} è attivo altrimenti è
  disattivo. Introdotta a partire dal kernel 3.4.
\itindend{child~reaper}


% TODO documentare PR_MPX_INIT e PR_MPX_RELEASE, vedi
% http://lwn.net/Articles/582712/ 

% TODO documentare PR_SET_MM_MAP aggiunta con il kernel 3.18, per impostare i
% parametri di base del layout dello spazio di indirizzi di un processo (area
% codice e dati, stack, brack pointer ecc. vedi
% http://git.kernel.org/linus/f606b77f1a9e362451aca8f81d8f36a3a112139e 

% TODO documentare ARCH_SET_CPUID e ARCH_GET_CPUID, introdotte con il kernel
% 4.12, vedi https://lwn.net/Articles/721182/

% TODO documentare PR_SPEC_DISABLE_NOEXEC in 5.1, vedi
% https://lwn.net/Articles/782511/ 

\label{sec:prctl_operation}
\end{basedescript}


\subsection{La funzione \func{ptrace}}
\label{sec:process_ptrace}

%Da fare

% TODO: trattare PTRACE_SEIZE, aggiunta con il kernel 3.1
% TODO: trattare PTRACE_O_EXITKILL, aggiunta con il kernel 3.8 (vedi
% http://lwn.net/Articles/529060/) 
% TODO: trattare PTRACE_GETSIGMASK e PTRACE_SETSIGMASK introdotte con il
% kernel 3.11
% TODO: trattare PTRACE_O_SUSPEND_SECCOMP, aggiunta con il kernel 4.3, vedi
% http://lwn.net/Articles/656675/ 

\subsection{La funzione \func{kcmp}}
\label{sec:process_kcmp}

% TODO: trattare kcmp aggiunta con il kernel 3.5, vedi
% https://lwn.net/Articles/478111/
% vedi man kcmp e man 2 open



\section{La gestione avanzata della creazione dei processi}
\label{sec:process_adv_creation}

In questa sezione tratteremo le funzionalità avanzate relative alla creazione
dei processi e del loro ambiente, sia per quanto riguarda l'utilizzo delle
stesse per la creazione dei \textit{thread} che per la gestione dei
\textit{namespace} che sono alla base dei cosiddetti \textit{container}.


\subsection{La \textit{system call} \func{clone}}
\label{sec:process_clone}

La funzione tradizionale con cui creare un nuovo processo in un sistema
Unix-like, come illustrato in sez.~\ref{sec:proc_fork}, è \func{fork}, ma con
l'introduzione del supporto del kernel per i \textit{thread}\unavref{ (vedi
  cap.~\ref{cha:threads})}, si è avuta la necessità di una interfaccia che
consentisse un maggiore controllo sulla modalità con cui vengono creati nuovi
processi, che poi è stata utilizzata anche per fornire supporto per le
tecnologie di virtualizzazione dei processi (i cosiddetti \textit{container})
su cui torneremo in sez.~\ref{sec:process_namespaces}.

Per questo l'interfaccia per la creazione di un nuovo processo è stata
delegata ad una nuova \textit{system call}, \funcm{sys\_clone}, che consente
di reimplementare anche la tradizionale \func{fork}. In realtà in questo caso
più che di nuovi processi si può parlare della creazioni di nuovi
``\textit{task}'' del kernel che possono assumere la veste sia di un processo
classico isolato dagli altri come quelli trattati finora, che di un
\textit{thread} in cui la memoria viene condivisa fra il processo chiamante ed
il nuovo processo creato, come quelli che vedremo in
sez.~\ref{sec:linux_thread}. Per evitare confusione fra \textit{thread} e
processi ordinari, abbiamo deciso di usare la nomenclatura \textit{task} per
indicare la unità di esecuzione generica messa a disposizione del kernel che
\texttt{sys\_clone} permette di creare.

La \textit{system call} richiede soltanto due argomenti: il
primo, \param{flags}, consente di controllare le modalità di creazione del
nuovo \textit{task}, il secondo, \param{child\_stack}, imposta l'indirizzo
dello \textit{stack} per il nuovo \textit{task}, e deve essere indicato quando
si intende creare un \textit{thread}. L'esecuzione del programma creato da
\func{sys\_clone} riprende, come per \func{fork}, da dopo l'esecuzione della
stessa.

La necessità di avere uno \textit{stack} alternativo c'è solo quando si
intende creare un \textit{thread}, in tal caso infatti il nuovo \textit{task}
vede esattamente la stessa memoria del \textit{task}
``\textsl{padre}'',\footnote{in questo caso per padre si intende semplicemente
  il \textit{task} che ha eseguito \func{sys\_clone} rispetto al \textit{task}
  da essa creato, senza nessuna delle implicazioni che il concetto ha per i
  processi.} e nella sua esecuzione alla prima chiamata di una funzione
andrebbe a scrivere sullo \textit{stack} usato anche dal padre (si ricordi
quanto visto in sez.~\ref{sec:proc_mem_layout} riguardo all'uso dello
\textit{stack}).

Per evitare di doversi garantire contro la evidente possibilità di
\textit{race condition} che questa situazione comporta (vedi
sez.~\ref{sec:proc_race_cond} per una spiegazione della problematica) è
necessario che il chiamante allochi preventivamente un'area di memoria.  In
genere lo si fa con una \func{malloc} che allochi un buffer che la funzione
imposterà come \textit{stack} del nuovo processo, avendo ovviamente cura di
non utilizzarlo direttamente nel processo chiamante.

In questo modo i due \textit{task} avranno degli \textit{stack} indipendenti e
non si dovranno affrontare problematiche di \textit{race condition}.  Si tenga
presente inoltre che in molte architetture di processore lo \textit{stack}
cresce verso il basso, pertanto in tal caso non si dovrà specificare
per \param{child\_stack} il puntatore restituito da \func{malloc}, ma un
puntatore alla fine del buffer da essa allocato.

Dato che tutto ciò è necessario solo per i \textit{thread} che condividono la
memoria, la \textit{system call}, a differenza della funzione di libreria che
vedremo a breve, consente anche di passare per \param{child\_stack} il valore
\val{NULL}, che non imposta un nuovo \textit{stack}. Se infatti si crea un
processo, questo ottiene un suo nuovo spazio degli indirizzi (è sottinteso
cioè che non si stia usando il flag \const{CLONE\_VM} che vedremo a breve) ed
in questo caso si applica la semantica del \textit{copy on write} illustrata
in sez.~\ref{sec:proc_fork}, per cui le pagine dello \textit{stack} verranno
automaticamente copiate come le altre e il nuovo processo avrà un suo
\textit{stack} totalmente indipendente da quello del padre.

Dato che l'uso principale della nuova \textit{system call} è quello relativo
alla creazione dei \textit{thread}, la \acr{glibc} definisce una funzione di
libreria con una sintassi diversa, orientata a questo scopo, e la
\textit{system call} resta accessibile solo se invocata esplicitamente come
visto in sez.~\ref{sec:proc_syscall}.\footnote{ed inoltre per questa
  \textit{system call} non è disponibile la chiamata veloce con
  \texttt{vsyscall}.} La funzione di libreria si chiama semplicemente
\funcd{clone} ed il suo prototipo è:

\begin{funcproto}{ 
\fhead{sched.h}
\fdecl{int clone(int (*fn)(void *), void *child\_stack, int flags, void *arg,
  ...  \\
\phantom{int clone(}/* pid\_t *ptid, struct user\_desc *tls, pid\_t *ctid */ )}
\fdesc{Crea un nuovo processo o \textit{thread}.} 
}
{La funzione ritorna il \textit{Thread ID} assegnato al nuovo processo in caso
  di successo e $-1$ per un errore, nel qual caso \var{errno} assumerà uno dei
  valori: 
\begin{errlist}
    \item[\errcode{EAGAIN}] sono già in esecuzione troppi processi.
    \item[\errcode{EINVAL}] si è usata una combinazione non valida di flag o
      un valore nullo per \param{child\_stack}.
    \item[\errcode{ENOMEM}] non c'è memoria sufficiente per creare una nuova
      \texttt{task\_struct} o per copiare le parti del contesto del chiamante
      necessarie al nuovo \textit{task}.
    \item[\errcode{EPERM}] non si hanno i privilegi di amministratore
      richiesti dai flag indicati.
\end{errlist}}
\end{funcproto}

% NOTE: una pagina con la descrizione degli argomenti:
% * http://www.lindevdoc.org/wiki/Clone 

La funzione prende come primo argomento \param{fn} il puntatore alla funzione
che verrà messa in esecuzione nel nuovo processo, che può avere un unico
argomento di tipo puntatore a \ctyp{void}, il cui valore viene passato dal
terzo argomento \param{arg}. Per quanto il precedente prototipo possa
intimidire nella sua espressione, in realtà l'uso è molto semplice basterà
definire una qualunque funzione \param{fn} che restituisce un intero ed ha
come argomento un puntatore a \ctyp{void}, e \code{fn(arg)} sarà eseguita in
un nuovo processo.

Il nuovo processo resterà in esecuzione fintanto che la funzione \param{fn}
non ritorna, o esegue \func{exit} o viene terminata da un segnale. Il valore
di ritorno della funzione (o quello specificato con \func{exit}) verrà
utilizzato come stato di uscita della funzione. I tre
argomenti \param{ptid}, \param{tls} e \param{ctid} sono opzionali e sono
presenti solo a partire dal kernel 2.6 e sono stati aggiunti come supporto per
le funzioni di gestione dei \textit{thread} (la \textit{Native Thread Posix
  Library}, vedi sez.~\ref{sec:linux_ntpl}) nella \acr{glibc}, essi vengono
utilizzati soltanto se si sono specificati rispettivamente i flag
\const{CLONE\_PARENT\_SETTID}, \const{CLONE\_SETTLS} e
\const{CLONE\_CHILD\_SETTID}. 

La funzione ritorna un l'identificatore del nuovo \textit{task}, denominato
\texttt{Thread ID} (da qui in avanti \ids{TID}) il cui significato è analogo
al \ids{PID} dei normali processi e che a questo corrisponde qualora si crei
un processo ordinario e non un \textit{thread}.

Il comportamento di \func{clone}, che si riflette sulle caratteristiche del
nuovo processo da essa creato, è controllato principalmente
dall'argomento \param{flags}, che deve essere specificato come maschera
binaria, ottenuta con un OR aritmetico di una delle costanti del seguente
elenco, che illustra quelle attualmente disponibili:\footnote{si fa
  riferimento al momento della stesura di questa sezione, cioè con il kernel
  3.2.}

\begin{basedescript}{\desclabelwidth{1.5 cm}\desclabelstyle{\nextlinelabel}}

\item[\constd{CLONE\_CHILD\_CLEARTID}] cancella il valore del \textit{thread
    ID} posto all'indirizzo dato dall'argomento \param{ctid}, eseguendo un
  riattivazione del \textit{futex} (vedi sez.~\ref{sec:xxx_futex}) a
  quell'indirizzo. Questo flag viene utilizzato dalla librerie di gestione dei
  \textit{thread} ed è presente dal kernel 2.5.49.

\item[\constd{CLONE\_CHILD\_SETTID}] scrive il \ids{TID} del \textit{thread}
  figlio all'indirizzo dato dall'argomento \param{ctid}. Questo flag viene
  utilizzato dalla librerie di gestione dei \textit{thread} ed è presente dal
  kernel 2.5.49.

\item[\constd{CLONE\_FILES}] se impostato il nuovo processo condividerà con il
  padre la \textit{file descriptor table} (vedi sez.~\ref{sec:file_fd}),
  questo significa che ogni \textit{file descriptor} aperto da un processo
  verrà visto anche dall'altro e che ogni chiusura o cambiamento dei
  \textit{file descriptor flag} di un \textit{file descriptor} verrà per
  entrambi.

  Se non viene impostato il processo figlio eredita una copia della
  \textit{file descriptor table} del padre e vale la semantica classica della
  gestione dei \textit{file descriptor}, che costituisce il comportamento
  ordinario di un sistema unix-like e che illustreremo in dettaglio in
  sez.~\ref{sec:file_shared_access}.

\item[\constd{CLONE\_FS}] se questo flag viene impostato il nuovo processo
  condividerà con il padre le informazioni relative all'albero dei file, ed in
  particolare avrà la stessa radice (vedi sez.~\ref{sec:file_chroot}), la
  stessa directory di lavoro (vedi sez.~\ref{sec:file_work_dir}) e la stessa
  \textit{umask} (sez.~\ref{sec:file_perm_management}). Una modifica di una
  qualunque di queste caratteristiche in un processo, avrà effetto anche
  sull'altro. Se assente il nuovo processo riceverà una copia delle precedenti
  informazioni, che saranno così indipendenti per i due processi, come avviene
  nel comportamento ordinario di un sistema unix-like.

\item[\constd{CLONE\_IO}] se questo flag viene impostato il nuovo processo
  condividerà con il padre il contesto dell'I/O, altrimenti, come avviene nel
  comportamento ordinario con una \func{fork} otterrà un suo contesto
  dell'I/O.

  Il contesto dell'I/O viene usato dagli \textit{scheduler} di I/O (visti in
  sez.~\ref{sec:io_priority}) e se questo è lo stesso per diversi processi
  questi vengono trattati come se fossero lo stesso, condividendo il tempo per
  l'accesso al disco, e possono interscambiarsi nell'accesso a disco. L'uso di
  questo flag consente, quando più \textit{thread} eseguono dell'I/O per conto
  dello stesso processo (ad esempio con le funzioni di I/O asincrono di
  sez.~\ref{sec:file_asyncronous_io}), migliori prestazioni.

%TODO : tutti i CLONE_NEW* attengono ai namespace, ed è meglio metterli nella
%relativa sezione da creare a parte

% \item[\constd{CLONE\_NEWIPC}] è uno dei flag ad uso dei \textit{container},
%   introdotto con il kernel 2.6.19. L'uso di questo flag crea per il nuovo
%   processo un nuovo \textit{namespace} per il sistema di IPC, sia per quello
%   di SysV (vedi sez.~\ref{sec:ipc_sysv}) che, dal kernel 2.6.30, per le code
%   di messaggi POSIX (vedi sez.~\ref{sec:ipc_posix_mq}); si applica cioè a
%   tutti quegli oggetti che non vegono identificati con un \textit{pathname}
%   sull'albero dei file.

%   L'uso di questo flag richiede privilegi di amministratore (più precisamente
%   la capacità \const{CAP\_SYS\_ADMIN}) e non può essere usato in combinazione
%   con \const{CLONE\_SYSVSEM}. 

% \item[\constd{CLONE\_NEWNET}]
% \item[\constd{CLONE\_NEWNS}]
% \item[\constd{CLONE\_NEWPID}]
% \item[\constd{CLONE\_NEWUTS}]


% TODO trattare CLONE_NEWCGROUP introdotto con il kernel 4.6, vedi
% http://lwn.net/Articles/680566/ 

\item[\constd{CLONE\_PARENT}]
\item[\constd{CLONE\_PARENT\_SETTID}]
\item[\constd{CLONE\_PID}]

% TODO trattare CLONE_PIDFD introdotto con il kernel 5.2, vedi
% https://lwn.net/Articles/787963/ e anche https://lwn.net/Articles/789023/
  
\item[\constd{CLONE\_PTRACE}] se questo flag viene impostato ed il processo
  chiamante viene tracciato (vedi sez.~\ref{sec:process_ptrace}) anche il
  figlio viene tracciato. 

\item[\constd{CLONE\_SETTLS}]
\item[\constd{CLONE\_SIGHAND}]
\item[\constd{CLONE\_STOPPED}]
\item[\constd{CLONE\_SYSVSEM}]
\item[\constd{CLONE\_THREAD}]

\item[\constd{CLONE\_UNTRACED}] se questo flag viene impostato un processo non
  può più forzare \const{CLONE\_PTRACE} su questo processo.

\item[\constd{CLONE\_VFORK}] se questo flag viene impostato il chiamante viene
  fermato fintato che il figlio appena creato non rilascia la sua memoria
  virtuale con una chiamata a \func{exec} o \func{exit}, viene quindi
  replicato il comportamento di \func{vfork}.

\item[\constd{CLONE\_VM}] se questo flag viene impostato il nuovo processo
  condividerà con il padre la stessa memoria virtuale, e le scritture in
  memoria fatte da uno qualunque dei processi saranno visibili dall'altro,
  così come ogni mappatura in memoria (vedi sez.~\ref{sec:file_memory_map}). 

  Se non viene impostato il processo figlio otterrà una copia dello spazio
  degli indirizzi e si otterrà il comportamento ordinario di un processo di un
  sistema unix-like creato con la funzione \func{fork}.
\end{basedescript}



\subsection{La gestione dei \textit{namespace}}
\label{sec:process_namespaces}

\itindbeg{namespace}
Come accennato all'inizio di sez.~\ref{sec:process_clone} oltre al controllo
delle caratteristiche dei processi usate per la creazione dei \textit{thread},
l'uso di \func{clone} consente, ad uso delle nuove funzionalità di
virtualizzazione dei processi, di creare nuovi ``\textit{namespace}'' per una
serie di proprietà generali (come l'elenco dei \ids{PID}, l'albero dei file, i
\textit{mount point}, la rete, il sistema di IPC, ecc.).

L'uso dei ``\textit{namespace}'' consente creare gruppi di processi che vedono
le suddette proprietà in maniera indipendente fra loro. I processi di ciascun
gruppo vengono così eseguiti come in una sorta di spazio separato da quello
degli altri gruppi, che costituisce poi quello che viene chiamato un
\textit{container}.

\itindend{namespace}


\itindbeg{container}

\itindend{container}


%TODO sezione separata sui namespace 

%TODO trattare unshare, vedi anche http://lwn.net/Articles/532748/

%TODO: trattare la funzione setns e i namespace file descriptors (vedi
% http://lwn.net/Articles/407495/) introdotti con il kernel 3.0, altre
% informazioni su setns qui: http://lwn.net/Articles/532748/
% http://lwn.net/Articles/531498/



\section{Funzionalità avanzate e specialistiche}
\label{sec:process_special}


% TODO: trattare userfaultfd, introdotta con il 4.23, vedi
% http://man7.org/linux/man-pages/man2/userfaultfd.2.html 

% TODO: trattare process_vm_readv/process_vm_writev introdotte con il kernel
% 3.2, vedi http://man7.org/linux/man-pages/man2/process_vm_readv.2.html e i
% precedenti tentativi https://lwn.net/Articles/405346/


\subsection{La gestione delle operazioni in virgola mobile}
\label{sec:process_fenv}

Da fare.

% TODO eccezioni ed arrotondamenti per la matematica in virgola mobile 
% consultare la manpage di fenv, math_error, fpclassify, matherr, isgreater,
% isnan, nan, INFINITY


\subsection{L'accesso alle porte di I/O}
\label{sec:process_io_port}

%
% TODO l'I/O sulle porte di I/O 
% consultare le manpage di ioperm, iopl e outb
% non c'entra nulla qui, va trovato un altro posto (altri meccanismi di I/O in
% fileintro ?)

Da fare


%\subsection{La gestione di architetture a nodi multipli}
%\label{sec:process_NUMA}

% TODO trattare i cpuset, che attiene anche a NUMA, e che possono essere usati
% per associare l'uso di gruppi di processori a gruppi di processi (vedi
% manpage omonima)
% TODO trattare getcpu, che attiene anche a NUMA, mettere qui anche
% sched_getcpu, che potrebbe essere indipendente ma richiama getcpu

%TODO trattare le funzionalità per il NUMA
% vedi man numa e, mbind, get_mempolicy, set_mempolicy, 
% le pagine di manuale relative
% vedere anche dove metterle...

% \subsection{La gestione dei moduli}
% \label{sec:kernel_modules}

% da fare

%TODO trattare init_module e finit_module (quest'ultima introdotta con il
%kernel 3.8)

%%%% Altre cose di cui non è chiara la collocazione:

%TODO trattare membarrier, introdotta con il kernel 4.3
% vedi http://lwn.net/Articles/369567/ http://lwn.net/Articles/369640/
% http://git.kernel.org/cgit/linux/kernel/git/torvalds/linux.git/commit/?id=5b25b13ab08f616efd566347d809b4ece54570d1 
% vedi anche l'ulteriore opzione "expedited" introdotta con il kernel 4.14
% (https://lwn.net/Articles/728795/) 



%%% Local Variables:
%%% mode: latex
%%% TeX-master: "gapil"
%%% End:

%  LocalWords:  system call namespace prctl IRIX kernel sys int option long
%  LocalWords:  unsigned arg errno EACCESS EBADF EBUSY EFAULT EINVAL ENXIO PR
%  LocalWords:  EOPNOTSUPP EPERM CAPBSET READ capability sez tab capabilities
%  LocalWords:  bounding CAP SETPCAP DUMPABLE dump suid sgid UID DISABLE GET
%  LocalWords:  ENDIAN endianness BIG big endian LITTLE little PPC PowerPC ia
%  LocalWords:  FPEMU NOPRINT SIGFPE FPEXC point exception FP EXC SW ENABLE
%  LocalWords:  OVF overflow UND underflow RES INV DISABLED NONRECOV ASYNC AO
%  LocalWords:  KEEPCAPS pag exec SECURE KEEP CAPS securebits LOCKED NAME NUL
%  LocalWords:  char PDEATHSIG SIGCHLD fork PTRACER PID tracer process ptrace
%  LocalWords:  Security Modules ANY Yama SECCOMP secure computing seccomp vm
%  LocalWords:  STRICT strict FILTER filter SIGKILL TIMING STATISTICAL TSC fn
%  LocalWords:  TIMESTAMP timestamp Stamp Counter SIGSEGV UNALIGN SIGBUS MCE
%  LocalWords:  KILL siginfo MCEERR memory failure early kill CLEAR child cap
%  LocalWords:  reaper SUBREAPER init value result thread like flags stack FS
%  LocalWords:  race condition malloc NULL copy write glibc vsyscall sched RT
%  LocalWords:  void pid ptid struct desc tls ctid EAGAIN ENOMEM exit Posix
%  LocalWords:  Library PARENT SETTID SETTLS TID CLEARTID futex FILES table
%  LocalWords:  descriptor umask dell'I scheduler SIGHAND STOPPED SYSVSEM IPC
%  LocalWords:  UNTRACED VFORK vfork mount filesystem LSM Mandatory Access fs
%  LocalWords:  Control DAC MAC SELinux Smack Tomoyo AppArmor Discrectionary
%  LocalWords:  permitted inheritable effective fig security ADMIN forced new
%  LocalWords:  allowed dall' bound MODULE nell' all' capset sendmail SETGID
%  LocalWords:  setuid orig IMMUTABLE MKNOD OVERRIDE SEARCH CHOWN FSETID LOCK
%  LocalWords:  FOWNER saved FIXUP NOROOT AUDIT BLOCK SUSPEND SETFCAP group
%  LocalWords:  socket domain locking mlock mlockall shmctl mmap OWNER LEASE
%  LocalWords:  lease immutable append only mknod BIND SERVICE BROADCAST RAW
%  LocalWords:  broadcast multicast PACKET CHROOT chroot NICE PACCT RAWIO TTY
%  LocalWords:  accounting ioperm iopl RESOURCE CONFIG hangup vhangup SYSLOG
%  LocalWords:  WAKE ALARM CLOCK BOOTTIME REALTIME sticky NOATIME fcntl swap
%  LocalWords:  multicasting dell'IPC SysV trusted IOPRIO CLASS IDLE lookup
%  LocalWords:  scheduling dcookie NEWNS unshare nice NUMA ioctl journaling
%  LocalWords:  ext capget header hdrp datap const ESRCH SOURCE undef version
%  LocalWords:  libcap lcap obj to text dup clear DIFFERS get ncap caps ssize
%  LocalWords:  argument length all setpcap from string name proc cat capgetp
%  LocalWords:  capsetp getcap read sigreturn sysctl protected hardlinks tmp
%  LocalWords:  dell' symlink symlinks pathname TOCTTOU of
